\documentclass{article}
\usepackage[french]{babel}
\usepackage[T1]{fontenc}
\usepackage{lmodern}
\usepackage{geometry}
\usepackage{amsmath}
\usepackage{enumitem}
\usepackage{titlesec}
\usepackage{xcolor}
\usepackage{tcolorbox}

% Configuration de la page
\geometry{a4paper, margin=2.5cm}
\setlength{\parindent}{0pt}
\setlength{\parskip}{1em}

% Style des sections
\titleformat{\section}[block]{\Large\bfseries\filcenter}{}{1em}{}
\titleformat{\subsection}[block]{\large\bfseries}{}{1em}{}

\begin{document}

% \maketitle % A décommenter si vous voulez un titre global

\section*{Introduction : La Gamme de Production comme Pilier du Système d'Information}

\subsection*{Définition : Qu'est-ce qu'une gamme de production ?}

La gamme de production constitue le document central et structurant de tout système d'information industriel. Elle est la traduction concrète du savoir-faire de l'entreprise en instructions opérationnelles. Formellement, elle se définit comme la \textbf{suite logique et ordonnée des phases, sous-phases et opérations} nécessaires à la transformation d'une matière première en un produit fini conforme aux spécifications.

Au sein de cette structure, on distingue :
\begin{itemize}
    \item \textbf{La Phase} : un ensemble de travaux réalisés sur un même poste de travail (machine, îlot d'assemblage).
    \item \textbf{La Sous-phase} : une séquence d'opérations effectuées sans modification du positionnement de la pièce.
    \item \textbf{L'Opération} : le plus petit élément de la gamme, correspondant à une action technique unique (ex : un passage d'outil, une insertion).
\end{itemize}

\subsection*{Rôle central : Le lien entre Conception, Industrialisation et Production}

La gamme n'est pas un document isolé ; elle est l'articulation fondamentale du triptyque CAO (Conception) – GPAO/ERP (Industrialisation/Gestion) – FAO (Exécution). Son rôle est stratégique car elle \textbf{impacte directement les trois leviers de performance de l'entreprise} :

\begin{tcolorbox}[colback=blue!5!white, colframe=blue!75!black, title=Impact de la gamme]
\begin{description}
    \item[Coûts] : Elle définit les ressources consommées (matières, outillage, main-d'œuvre, temps machine). Une gamme optimisée est le premier facteur de maîtrise des coûts de revient.
    \item[Délais] : Le séquencement des opérations, le calcul des temps unitaires et des temps de préparation conditionnent directement les délais de fabrication et la réactivité face à la demande.
    \item[Qualité] : Elle intègre les procédures de contrôle, les conditions de coupe, les modes opératoires qui garantissent la conformité et la reproductibilité du produit.
\end{description}
\end{tcolorbox}

\subsection*{Le défi multi-classe : Pourquoi une approche unique ne fonctionne pas ?}

Si la gamme est un concept universel, sa logique de construction et de gestion diffère radicalement selon le contexte de production. Face à la diversité des marchés (grande série, moyenne série, unitaire), une approche unique et rigide est vouée à l'échec. Le système d'information doit s'adapter à ces trois classes de problématiques :

\begin{enumerate}[label=\textbf{Classe \Alph*}, leftmargin=*]
    \item \textbf{Produits standards, grande série : L'optimisation du flux}
    \begin{itemize}
        \item \textit{Objectif} : Maximiser la productivité et réduire les temps de cycle.
        \item \textit{Logique} : Standardisation, répétitivité, recherche du temps de cycle minimal. Les gammes sont figées et très détaillées. L'enjeu est l'équilibrage des lignes et l'ordonnancement pour un flux continu (ex : industrie automobile, électroménager).
    \end{itemize}
    
    \item \textbf{Produits à variantes, moyenne série : La modularité}
    \begin{itemize}
        \item \textit{Objectif} : Offrir de la diversité tout en maîtrisant la complexité industrielle.
        \item \textit{Logique} : Conception modulaire. Les gammes doivent gérer les options et les variantes de produits. L'approche est celle de la \textbf{nomenclature générique} et des \textbf{gammes de remplacement}. L'enjeu est de trouver l'équilibre entre standardisation (des sous-ensembles) et personnalisation (des options).
    \end{itemize}
    
    \item \textbf{Produits spécifiques, petite série / unitaire : La flexibilité}
    \begin{itemize}
        \item \textit{Objectif} : Répondre à une demande unique et personnalisée dans des délais contraints.
        \item \textit{Logique} : Approche par projet. Les gammes sont créées \textit{ad-hoc} ou par adaptation d'un existant. L'enjeu n'est plus l'optimisation marginale du temps de cycle, mais la \textbf{capacité à fiabiliser un processus inédit}, à gérer les risques techniques et à garantir une qualité irréprochable dès la première pièce (ex : aéronautique, outillage lourd).
    \end{itemize}
\end{enumerate}

En conclusion, définir une gamme, c'est choisir une stratégie industrielle. Le système d'information qui la porte doit donc être capable d'incarner ces trois logiques : de l'optimisation continue des flux à la gestion flexible de projets uniques, en passant par la configuration modulaire de produits complexes. C'est à cette condition qu'il peut véritablement être considéré comme un pilier stratégique et non comme un simple outil de saisie.
% ============================================
% PARTIE 1 : PRÉREQUIS ET FONDAMENTAUX
% ============================================

\section*{Partie 1 : Prérequis et Fondamentaux}
\subsection*{Ce qu'il faut savoir avant de définir une gamme}

Avant toute définition de gamme, une analyse rigoureuse du produit et de son environnement de fabrication est indispensable. Cette partie pose les fondements nécessaires à la construction d'une gamme cohérente et performante.

% ============================================
% CHAPITRE 1 : ANALYSE DU PRODUIT (LE "QUOI")
% ============================================

\section{Analyse du Produit (Le "Quoi")}

\subsection{Décomposition fonctionnelle}

La décomposition fonctionnelle consiste à analyser le produit non pas comme un objet géométrique, mais comme un ensemble de \textbf{fonctions de service} qu'il doit remplir. Cette approche, fondamentale en ingénierie concourante, permet de :

\begin{itemize}
    \item \textbf{Identifier les fonctions principales} : celles pour lesquelles le produit est conçu (ex : supporter une charge, transmettre un mouvement).
    \item \textbf{Caractériser les fonctions de contrainte} : celles liées à l'environnement (ex : résistance à la corrosion, encombrement).
    \item \textbf{Associer chaque fonction à des zones usinées} : cette mise en correspondance oriente les choix technologiques et les tolérances.
\end{itemize}

\begin{tcolorbox}[colback=gray!5!white, colframe=gray!75!black, title=Méthodologie de décomposition]
\begin{enumerate}
    \item \textbf{Lister les fonctions} : "Le produit doit..."
    \item \textbf{Hiérarchiser} : Fonctions principales / fonctions secondaires / contraintes
    \item \textbf{Critères d'appréciation} : Niveau de performance attendu (ex : effort max, précision dimensionnelle)
    \item \textbf{Traduction géométrique} : Quelles surfaces participent à quelle fonction ?
\end{enumerate}
\end{tcolorbox}

\subsection{Construction de la Nomenclature (BOM)}

La \textbf{Bill of Materials (BOM)} ou nomenclature est la structure arborescente qui décrit la composition du produit. Elle constitue le squelette de toute gestion de production.

\subsubsection{Types de nomenclatures}

\begin{tabular}{|p{0.3\textwidth}|p{0.65\textwidth}|}
\hline
\textbf{Type de nomenclature} & \textbf{Caractéristiques} \\
\hline
\textbf{Nomenclature d'étude} & Structure conceptuelle, souvent issue de la CAO, indépendante des contraintes de fabrication. \\
\hline
\textbf{Nomenclature de production} & Structure industrialisée, intégrant les contraintes de fabrication (sous-ensembles, gammes associées). \\
\hline
\textbf{Nomenclature multi-niveaux} & Représentation hiérarchique où chaque composant peut être lui-même décomposé. \\
\hline
\textbf{Nomenclature mono-niveau} & Vue "père-fils" immédiate, utilisée pour le calcul des besoins. \\
\hline
\textbf{Nomenclature de planification} & Regroupement d'articles en familles pour la planification agrégée (Plan Industriel et Commercial). \\
\hline
\end{tabular}

\subsubsection{Principes de structuration}

\begin{center}
\begin{tcolorbox}[colback=white, colframe=blue!50!black, width=0.8\textwidth]
\textbf{Représentation arborescente de la Brouette :}
\vspace{0.5em}

\begin{tabular}{ll}
Niveau 1 & \textbf{Brouette} \\
& $\quad$ Niveau 2 : Structure + Support \\
& $\quad$ Niveau 2 : Benne \\
& $\quad$ Niveau 2 : Ensemble portant \\
& $\qquad$ Niveau 3 : Roue \\
& $\qquad$ Niveau 3 : Axe \\
& $\qquad$ Niveau 3 : Manchons \\
& $\qquad$ Niveau 3 : Boulons \\
& $\quad$ Niveau 2 : Peinture \\
\end{tabular}
\end{tcolorbox}
\end{center}

\subsection{Classification des Articles (Acheté, Fabriqué, Générique, Spécifique)}

La classification des articles est un préalable essentiel qui détermine leur mode de gestion et leur place dans la gamme.

\begin{description}[font=\bfseries\color{blue!70!black}]
    \item[Article Acheté] Produit approvisionné chez un fournisseur externe. Il constitue généralement une feuille de la nomenclature et peut être soumis à des contrôles qualité à réception.
    
    \item[Article Fabriqué] Sous-ensemble ou produit final réalisé dans l'entreprise (ou chez un sous-traitant). Il est associé à une gamme de fabrication spécifique.
    
    \item[Article Générique] Article existant sous de multiples variantes. Sa configuration est définie au moment de la commande client. Exemple : une perceuse électrique déclinée en puissance et en couleur.
    
    \item[Article Spécifique] Article réalisé selon des spécifications client uniques pour un projet donné. Il peut être dérivé d'un article standard ou générique. Sa gestion est généralement rattachée à un projet PCS (Projet Client Spécifique).
\end{description}

\begin{tcolorbox}[colback=yellow!10!white, colframe=yellow!50!black, title=Conseils pratiques]
\begin{itemize}
    \item Un article "Fabriqué" nécessite toujours une gamme et une nomenclature.
    \item Un article "Acheté" peut être lié à une fiche fournisseur et un délai d'approvisionnement.
    \item La distinction entre "Générique" et "Spécifique" est cruciale pour le paramétrage du système : un article générique peut être configuré à la commande, un article spécifique est unique et rattaché à un projet.
\end{itemize}
\end{tcolorbox}

\subsection{Identification des données techniques par classe de produit}

Chaque classe de produit impose des données techniques spécifiques à formaliser avant la construction de la gamme.

\subsubsection{Produits de série (Classe A \& B)}

\begin{itemize}
    \item \textbf{Données de base} : Matière, dimensions brutes et finies, poids.
    \item \textbf{Données de fabrication} : Temps unitaires standards, temps de préparation, cadences.
    \item \textbf{Données d'outillage} : Outils dédiés, porte-outils, conditions de coupe optimisées.
    \item \textbf{Données logistiques} : Taille de lot économique, stock de sécurité, seuil de réapprovisionnement.
\end{itemize}

\subsubsection{Produits à la commande (Classe C - Projet)}

\begin{itemize}
    \item \textbf{Dossier technique projet} : Plan client, spécifications particulières, référentiel qualité.
    \item \textbf{Données de risques} : Identification des opérations critiques, analyse des modes de défaillance.
    \item \textbf{Traçabilité} : Numéro d'évolution, rattachement commande client, lots matière.
    \item \textbf{Phasage projet} : Jalons, dates clés, activités préalables (ex : contrôle métallurgique avant usinage).
\end{itemize}

\begin{tcolorbox}[colback=green!5!white, colframe=green!50!black, title=Synthèse : Ce qu'il faut savoir]
\begin{center}
\textbf{Avant de définir une gamme, le méthodiste doit disposer :}
\end{center}
\begin{enumerate}
    \item De la \textbf{nomenclature} (structure du produit).
    \item Du \textbf{classement} de chaque article (acheté/fabriqué/générique/spécifique).
    \item Des \textbf{spécifications géométriques} et fonctionnelles.
    \item Des \textbf{ressources disponibles} (machines, outillages, compétences).
    \item Des \textbf{contraintes} (série, délais, coûts cibles, exigences qualité).
\end{enumerate}
\end{tcolorbox}
% ============================================
% CHAPITRE 2 : ANALYSE DES MOYENS DE PRODUCTION (LE "AVEC QUOI")
% ============================================

\section{Analyse des Moyens de Production (Le "Avec Quoi")}

L'analyse des moyens de production constitue le second pilier fondamental de la préparation de gamme. Elle répond à la question centrale : \textbf{quels sont les ressources et équipements disponibles pour réaliser la fabrication ?} Cette analyse conditionne à la fois la faisabilité technique et la performance économique des gammes.

\subsection{Inventaire des Ressources (Machines, Postes de charge, Outillages)}

\subsubsection{Les Machines et Équipements de Production}

L'inventaire des machines doit être exhaustif et documenté. Chaque machine est caractérisée par :

\begin{itemize}
    \item \textbf{Identification} : Code machine, désignation, localisation
    \item \textbf{Caractéristiques techniques} : 
    \begin{itemize}
        \item Puissance moteur (kW) et couple disponible (Nm)
        \item Plage de vitesses (min et max, tr/min)
        \item Nombre d'axes (3, 4, 5 axes, axes positionnés ou continus)
        \item Course des axes (X, Y, Z)
        \item Capacité d'outillage (magasin, changeur automatique)
        \item Accessoires disponibles (têtes à renvoi d'angle, tables rotatives)
    \end{itemize}
    \item \textbf{Temps de changement} : 
    \begin{itemize}
        \item Temps de changement d'outil unitaire
        \item Temps de changement d'accessoire
        \item Temps de changement d'orientation
    \end{itemize}
\end{itemize}

\begin{tcolorbox}[colback=blue!5!white, colframe=blue!75!black, title=Exemple : Fiche machine pour atelier d'usinage]
\begin{tabular}{|p{0.35\textwidth}|p{0.6\textwidth}|}
\hline
\textbf{Caractéristique} & \textbf{Valeur} \\
\hline
Référence & Fraiseuse 3 axes + tête à renvoi d'angle \\
\hline
Puissance broche & 100 kW \\
\hline
Vitesse max broche & 1500 tr/min \\
\hline
Courses (X/Y/Z) & 2000 / 1000 / 800 mm \\
\hline
Changeur d'outils & 60 positions, temps unitaire 0,5 min \\
\hline
Tête à renvoi d'angle & Orientation 0-90-180-270°, temps changement 5 min \\
\hline
\end{tabular}
\end{tcolorbox}

\subsubsection{Les Postes de Charge}

Le poste de charge est l'unité élémentaire de planification de la production. Il peut être :

\begin{itemize}
    \item \textbf{Un centre de charge unique} : une machine spécifique
    \item \textbf{Un groupe de machines} : plusieurs machines aux caractéristiques similaires
    \item \textbf{Un poste manuel} : îlot d'assemblage ou de contrôle
\end{itemize}

\begin{tcolorbox}[colback=gray!5!white, colframe=gray!75!black, title=Caractérisation d'un poste de charge]
\begin{description}
    \item[Code Poste] Identifiant unique dans le système
    \item[Type] Machine, cellule flexible, poste manuel
    \item[Capacité] Nombre de personnes/machines affectées
    \item[Calendrier] Jours ouvrés, équipes, pauses
    \item[Compétences] Qualifications requises des opérateurs
    \item[Taux horaire] Coût machine et main-d'œuvre associés
\end{description}
\end{tcolorbox}

\subsubsection{L'Outillage}

L'outillage représente un actif stratégique souvent sous-estimé. Son inventaire doit inclure :

\begin{itemize}
    \item \textbf{Outils coupants} :
    \begin{itemize}
        \item Type : Fraise, foret, alésoir, taraud
        \item Géométrie : Diamètre, longueur, rayon de bec, nombre de dents
        \item Matériau : Carbure, HSS, Céramique
        \item Données de coupe : Vitesse de coupe (Vc), avance à la dent (Fz)
        \item Durée de vie : Temps d'usinage avant changement
        \item Coût : Prix unitaire, coût par arête
    \end{itemize}
    \item \textbf{Outillage de montage} :
    \begin{itemize}
        \item Bridage, étaux, plaques magnétiques
        \item Montages spécifiques par famille de produits
        \item Système de palettisation
    \end{itemize}
    \item \textbf{Outils de contrôle} :
    \begin{itemize}
        \item Palpeurs, capteurs
        \item Instruments de mesure (pied à coulisse, micromètre, MMT)
    \end{itemize}
\end{itemize}

\subsection{Définition des Calendriers et des Capacités}

La définition précise des calendriers et des capacités est essentielle pour la planification réaliste des ordres de fabrication.

\subsubsection{Les Types de Calendriers}

\begin{tabular}{|p{0.3\textwidth}|p{0.65\textwidth}|}
\hline
\textbf{Type de calendrier} & \textbf{Usage} \\
\hline
\textbf{Calendrier société} & Jours fériés, fermeture annuelle, week-ends \\
\hline
\textbf{Calendrier poste de charge} & Disponibilité spécifique d'un centre de charge \\
\hline
\textbf{Calendrier machine} & Maintenance préventive, indisponibilités techniques \\
\hline
\textbf{Calendrier équipe} & Roulement, horaires, pauses \\
\hline
\end{tabular}

\subsubsection{Définition des Capacités}

\begin{description}
    \item[Capacité théorique] Nombre d'heures disponibles sur la base d'un fonctionnement continu.
    \[
    C_{théo} = N_{jours} \times N_{heures/jour} \times N_{postes} \times N_{machines}
    \]
    
    \item[Capacité pratique] Capacité théorique diminuée des temps d'arrêts programmés (maintenance, changements de série, pauses).
    \[
    C_{pratique} = C_{théo} - T_{arrêts programmés}
    \]
    
    \item[Capacité utilisable] Capacité pratique affectée d'un coefficient d'efficacité (rendement opérationnel).
    \[
    C_{utilisable} = C_{pratique} \times \eta
    \]
    Où $\eta$ (efficacité) intègre les aléas, la fatigue, les temps de réglage non anticipés.
\end{description}

\begin{tcolorbox}[colback=yellow!10!white, colframe=yellow!50!black, title=Bonnes pratiques]
\begin{itemize}
    \item \textbf{Définir des horizons de planification} : horizon figé (ordres confirmés) et horizon glissant (prévisions).
    \item \textbf{Paramétrer des seuils de capacité} : alerte lorsque le taux d'utilisation dépasse un seuil critique (ex : 85\%).
    \item \textbf{Identifier les goulets d'étranglement} : ressources critiques dont la capacité limite le flux de production.
\end{itemize}
\end{tcolorbox}

\subsection{Établissement des Centres de Coût et des Taux Horaires}

La maîtrise des coûts passe par une structure claire des centres de coût et des taux horaires associés.

\subsubsection{Définition des Centres de Coût}

Un centre de coût est une entité organisationnelle (machine, atelier, service) pour laquelle on collecte et analyse les charges. On distingue :

\begin{itemize}
    \item \textbf{Centres de coût directs} : directement liés à la production (machines, postes d'assemblage)
    \item \textbf{Centres de coût indirects} : support à la production (maintenance, méthodes, qualité)
    \item \textbf{Centres de coût de structure} : administration, direction
\end{itemize}

\subsubsection{Calcul du Taux Horaire Machine}

Le taux horaire machine (ou \textbf{taux opératoire}) représente le coût d'utilisation d'un centre de charge par unité de temps.

\begin{tcolorbox}[colback=green!5!white, colframe=green!50!black, title=Formule de calcul du taux horaire]
\[
THM = \frac{C_{fixes} + C_{variables} + C_{indirects}}{H_{facturables}}
\]
\vspace{0.5em}
Avec :
\begin{itemize}
    \item $C_{fixes}$ : Amortissement, locaux, assurances, entretien
    \item $C_{variables}$ : Énergie, consommables, outillage
    \item $C_{indirects}$ : Part des services supports
    \item $H_{facturables}$ : Nombre d'heures facturables par an
\end{itemize}
\end{tcolorbox}

\subsubsection{Composantes du Taux Horaire}

\begin{tabular}{|p{0.35\textwidth}|p{0.6\textwidth}|}
\hline
\textbf{Composante} & \textbf{Description} \\
\hline
\textbf{Amortissement} & Valeur de la machine répartie sur sa durée de vie \\
\hline
\textbf{Maintenance} & Coûts de maintenance préventive et curative \\
\hline
\textbf{Énergie} & Consommation électrique, pneumatique, fluides \\
\hline
\textbf{Outillage} & Coûts des outils de coupe, porte-outils, montages \\
\hline
\textbf{Main-d'œuvre} & Salaire et charges de l'opérateur affecté \\
\hline
\textbf{Frais généraux} & Encadrement, locaux, assurance, taxes \\
\hline
\end{tabular}

\subsubsection{Taux horaire main-d'œuvre}

Pour les postes manuels, le taux horaire est calculé à partir :

\[
THMO = \frac{Salaire_{annuel} + Charges_{sociales} + Frais_{structure}}{Heures_{travaillées}}
\]

\begin{tcolorbox}[colback=blue!5!white, colframe=blue!75!black, title=Exemple : Taux horaire machine de fraisage 3 axes]
\begin{tabular}{|l|r|}
\hline
\textbf{Composante} & \textbf{Coût annuel (\textdollar)} \\
\hline
Amortissement (10 ans) & 35 000 \\
Maintenance & 8 000 \\
Énergie & 5 000 \\
Outillage & 12 000 \\
Frais généraux (30\%) & 18 000 \\
\hline
\textbf{Total coûts annuels} & \textbf{78 000} \\
\hline
Heures facturables (1800 h/an) &  \\
\hline
\textbf{Taux horaire} & \textbf{43,33 \$/h} \\
\hline
\end{tabular}
\end{tcolorbox}

\subsubsection{Intégration dans le calcul du coût de revient}

Le taux horaire est utilisé pour valoriser les temps de fabrication dans le calcul du coût de revient :

\[
Coût_{fabrication} = \sum_{opérations} (T_{usinage} \times THM) + \sum_{opérations} (T_{manuel} \times THMO) + Coût_{matières} + Coût_{outils}
\]

\begin{tcolorbox}[colback=gray!5!white, colframe=gray!75!black, title=Synthèse : Analyse des Moyens de Production]
\begin{center}
\textbf{Pour une analyse exhaustive, le méthodiste doit constituer :}
\end{center}
\begin{enumerate}
    \item Un \textbf{inventaire des machines} avec leurs caractéristiques techniques et temps de changement
    \item Une \textbf{description des postes de charge} et de leurs calendriers de disponibilité
    \item Un \textbf{catalogue des outillages} avec leurs données technico-économiques
    \item Une \textbf{grille des taux horaires} par centre de coût
    \item Une \textbf{matrice des compétences} associant opérateurs et postes
\end{enumerate}
\textit{Cette base de données structurée est le socle indispensable pour toute optimisation de gamme.}
\end{tcolorbox}
% ============================================
% CHAPITRE 3 : ANALYSE DES CONTRAINTES DE FABRICATION (LE "COMMENT")
% ============================================

\section{Analyse des Contraintes de Fabrication (Le "Comment")}

L'analyse des contraintes de fabrication constitue le troisième pilier fondamental. Elle répond à la question : \textbf{comment réaliser la fabrication en respectant les spécifications du produit et les limites des moyens disponibles ?} Cette analyse conditionne la faisabilité, la qualité et la robustesse des gammes.

\subsection{Identification des contraintes géométriques}

Les contraintes géométriques sont directement issues du dessin de définition et du modèle CAO. Elles conditionnent l'accessibilité des outils, les stratégies d'usinage et le choix des montages.

\subsubsection{Contraintes d'accessibilité}

\begin{tcolorbox}[colback=blue!5!white, colframe=blue!75!black, title=Contraintes d'accessibilité - Définitions]
\begin{description}
    \item[Accessibilité directe] L'outil peut atteindre la surface sans interférence avec la pièce, le montage ou la machine.
    \item[Accessibilité indirecte] Nécessité d'un outil de forme particulière (fraise longue, tête à renvoi d'angle) ou d'un repositionnement de la pièce.
    \item[Zone inaccessible] Surface ne pouvant être usinée avec les moyens disponibles, nécessitant une adaptation du design ou un procédé alternatif.
\end{description}
\end{tcolorbox}

Les paramètres géométriques clés à identifier pour chaque zone à usiner sont :

\begin{tabular}{|p{0.35\textwidth}|p{0.6\textwidth}|}
\hline
\textbf{Paramètre} & \textbf{Impact sur la gamme} \\
\hline
\textbf{Rayon concave minimum} & Conditionne le diamètre maximal de l'outil de finition ($R_{outil} \leq R_{concave}$). \\
\hline
\textbf{Rayon outil maximal} & Limite le diamètre d'outil pouvant pénétrer dans la cavité. \\
\hline
\textbf{Longueur minimale d'outil} & Définit la longueur nécessaire pour atteindre toute la zone sans collision. \\
\hline
\textbf{Profondeur de cavité} & Impacte le choix de l'outil (longueur, diamètre) et le nombre de passes axiales. \\
\hline
\textbf{Épaisseur de matière à enlever} & Détermine le nombre de passes radiales et la stratégie d'ébauche. \\
\hline
\textbf{Variation d'engagement radial} & Influence le risque de vibrations et la stabilité de coupe. \\
\hline
\end{tabular}

\subsubsection{Contraintes dimensionnelles et tolérances}

\begin{itemize}
    \item \textbf{Tolérances dimensionnelles} : Impactent le choix entre ébauche, demi-finition, finition et super-finition.
    \item \textbf{Tolérances géométriques} : 
    \begin{itemize}
        \item Planéité, perpendicularité, coaxialité : imposent des stratégies d'usinage en un même posage.
        \item Positionnement : peut nécessiter des opérations de contrôle intermédiaire.
    \end{itemize}
    \item \textbf{État de surface} : 
    \begin{itemize}
        \item $Ra$ (rugosité moyenne) : conditionne l'avance à la dent et la stratégie de finition.
        \item $Rz$ (hauteur de crêtes) : limite le ressaut maximal admissible entre passes.
    \end{itemize}
\end{itemize}

\begin{tcolorbox}[colback=yellow!10!white, colframe=yellow!50!black, title=Exemple : Tolérances et stratégies d'usinage]
\begin{center}
\begin{tabular}{|c|c|}
\hline
\textbf{Tolérance / État de surface} & \textbf{Stratégie recommandée} \\
\hline
$Ra \leq 0.8 \ \mu m$ & Super-finition, outils à plaquettes wiper \\
\hline
$Ra \leq 1.6 \ \mu m$ & Finition, avance réduite \\
\hline
$Ra \leq 3.2 \ \mu m$ & Demi-finition \\
\hline
$Ra \geq 6.3 \ \mu m$ & Ébauche uniquement \\
\hline
\end{tabular}
\end{center}
\end{tcolorbox}

\subsection{Identification des contraintes technologiques}

Les contraintes technologiques proviennent de l'interaction entre le matériau usiné, l'outil et la machine.

\subsubsection{Contraintes liées au matériau}

\begin{description}
    \item[Matériau usiné] Alliage de titane, acier inoxydable, aluminium, etc. Chaque matériau possède ses propres caractéristiques d'usinabilité.
    \item[Force de coupe spécifique $K_c$] Détermine l'effort de coupe et donc la puissance nécessaire.
    \item[Module d'Young $E$] Impacte la flexion de l'outil et donc les risques de vibrations.
    \item[Conductivité thermique] Influence l'évacuation de la chaleur et l'usure de l'outil.
\end{description}

\begin{tcolorbox}[colback=green!5!white, colframe=green!50!black, title=Exemple : Usinage d'alliage de titane]
\begin{itemize}
    \item \textbf{Caractéristiques} : Haute résistance mécanique, faible conductivité thermique, forte affinité chimique avec les outils.
    \item \textbf{Conséquences} :
    \begin{itemize}
        \item Vitesses de coupe réduites ($V_c \approx 30-60 \ m/min$)
        \item Avances à la dent modérées ($F_z \approx 0.05-0.15 \ mm/dent$)
        \item Outils en carbure avec revêtement spécifique
        \item Utilisation intensive de lubrification
        \item Durée de vie outil limitée
    \end{itemize}
\end{itemize}
\end{tcolorbox}

\subsubsection{Contraintes liées à l'outil}

\begin{itemize}
    \item \textbf{Limites technologiques} :
    \begin{itemize}
        \item Section de copeau maximale admissible $A_{Smax}$
        \item Profondeur de passe axiale maximale $Ap_{max}$
        \item Engagement radial minimal $Ae_{min}$ et maximal $Ae_{max}$
        \item Avance à la dent minimale $F_{zmin}$ et maximale $F_{zmax}$
        \item Durée de vie outil $T_{outil}$ (heures de coupe)
    \end{itemize}
    \item \textbf{Contraintes de rupture} : Effort de coupe limite avant rupture de l'outil
    \item \textbf{Contraintes de flexion} : Flèche maximale admissible pour respecter les tolérances géométriques
\end{itemize}

\subsubsection{Contraintes liées à la machine}

\begin{itemize}
    \item \textbf{Contraintes de puissance} :
    \[
    P_{coupe} = \frac{V_c \times Ap \times Ae \times K_c}{60 \times 1000} \leq P_{broche}
    \]
    
    \item \textbf{Contraintes de couple} :
    \[
    C_{coupe} = \frac{P_{coupe} \times 60}{2\pi \times N} \leq C_{broche}
    \]
    
    \item \textbf{Contraintes cinématiques} :
    \begin{itemize}
        \item Vitesse de broche $N$ dans la plage $[N_{min}, N_{max}]$
        \item Vitesse d'avance $V_f$ dans la plage de la machine
    \end{itemize}
    
    \item \textbf{Contraintes d'encombrement} :
    \begin{itemize}
        \item Courses axes $X$, $Y$, $Z$ suffisantes pour la pièce
        \item Espace magasin outil pour tous les outils de la gamme
    \end{itemize}
\end{itemize}

\subsection{Identification des contraintes de séquencement (antériorité)}

Les contraintes de séquencement définissent l'ordre dans lequel les opérations doivent être réalisées pour garantir la faisabilité et la qualité.

\subsubsection{Contraintes d'antériorité entre zones}

Certaines zones doivent impérativement être usinées avant d'autres pour des raisons :

\begin{itemize}
    \item \textbf{D'accessibilité} : Une zone peut masquer l'accès à une autre.
    \item \textbf{De tenue de pièce} : L'usinage préalable d'une zone peut fragiliser la pièce.
    \item \textbf{De référencement} : Une zone sert de référence pour le positionnement de la pièce.
    \item \textbf{D'état de contrainte} : L'usinage peut libérer des contraintes internes.
\end{itemize}

\begin{tcolorbox}[colback=blue!5!white, colframe=blue!75!black, title=Formalisation des contraintes d'antériorité]
Soit $Zone_i$ et $Zone_j$ deux zones à usiner.
\[
Zone_i \prec Zone_j \ \text{(Zone i doit être usinée avant Zone j)}
\]
L'ensemble $Ant_{Zone_i} = \{ Zone_j \ | \ Zone_j \prec Zone_i \}$ représente les zones antérieures obligatoires.

\textbf{Contrainte générale :}
\[
\forall Zone_i, \ \forall n \in \{1,\dots,Nb_{opm}\} \text{ tel que } Zone_{opm,n} = Zone_i,
\]
\[
\exists j \in \{1,\dots,n-1\} \text{ tel que } Zone_{opm,j} \in Ant_{Zone_i}
\]
\end{tcolorbox}

\subsubsection{Contraintes d'antériorités entre opérations d'une même zone}

Pour une même zone, l'ordre logique des opérations est imposé par le principe de progression de l'usinage :

\begin{center}
\begin{tcolorbox}[colback=gray!5!white, colframe=gray!75!black, width=0.8\textwidth]
\textbf{Hiérarchie des opérations :}
\vspace{0.5em}

\[
\text{Ébauche} \rightarrow \text{Reprise} \rightarrow \text{Demi-finition} \rightarrow \text{Finition} \rightarrow \text{Super-finition}
\]
\end{tcolorbox}
\end{center}

\begin{itemize}
    \item \textbf{Ébauche} : Enlèvement de matière massif, tolérances larges.
    \item \textbf{Reprise} : Correction des défauts d'ébauche, préparation pour la demi-finition.
    \item \textbf{Demi-finition} : Approche des dimensions finales, préparation de l'état de surface.
    \item \textbf{Finition} : Atteinte des cotes et états de surface spécifiés.
    \item \textbf{Super-finition} : Opération de très haute précision (ex : rodage, polissage).
\end{itemize}

\subsubsection{Contraintes de regroupement et d'optimisation}

\begin{description}
    \item[Regroupement par outil] Opérations réalisables avec le même outil peuvent être regroupées pour réduire les changements d'outil.
    
    \item[Regroupement par orientation] Opérations dans la même orientation peuvent être séquencées pour minimiser les changements d'axe.
    
    \item[Regroupement par accessoire] Opérations nécessitant le même accessoire peuvent être regroupées pour réduire les changements d'accessoire.
    
    \item[Contrainte de préparation] Certaines opérations nécessitent une préparation spécifique (ex : contrôle de l'outil, réglage machine).
\end{description}

\subsubsection{Contraintes liées à la qualité et au contrôle}

\begin{itemize}
    \item \textbf{Contrôle intermédiaire} : Après certaines opérations critiques, un contrôle peut être obligatoire avant la suite.
    \item \textbf{Traceabilité} : Les opérations doivent être enregistrées dans l'ordre pour garantir la traçabilité.
    \item \textbf{Points de comptage} : Opérations où la quantité achevée doit être déclarée explicitement pour le suivi des rebuts.
\end{itemize}

\begin{tcolorbox}[colback=yellow!10!white, colframe=yellow!50!black, title=Exemple : Graphe d'antériorité pour une pièce de structure]
\begin{center}
\begin{tabular}{c}
Zone extérieure (référence) \\
$\downarrow$ \\
Zone de perçage (pré-perçage avant ébauche) \\
$\downarrow$ \\
Ébauche zone intérieure \\
$\downarrow$ \\
Demi-finition zone intérieure \\
$\downarrow$ \\
Finition zone intérieure \\
$\downarrow$ \\
Finition zone extérieure (reprise de cote) \\
$\downarrow$ \\
Perçage final \\
\end{tabular}
\end{center}
\textit{Remarque : Les opérations de finition sont réalisées en dernier pour éviter toute déformation ultérieure.}
\end{tcolorbox}

\subsubsection{Formalisation mathématique des contraintes de séquencement}

Pour une gamme $m$ composée de $Nb_{opm}$ opérations $Opm_{m,n}$ :

\begin{itemize}
    \item \textbf{Contrainte d'antérioritè entre zones} :
    \[
    \forall n \in \{1,\dots,Nb_{opm}-1\}, \ \forall Zone \in Ant_{Zone_{opm,n}}, \ \exists j \in \{1,\dots,n-1\} \text{ tel que } Zone_{opm,j} = Zone
    \]
    
    \item \textbf{Contrainte d'antériorité entre types d'opération pour une même zone} :
    \[
    \forall n \in \{1,\dots,Nb_{opm}-1\} \text{ avec } Zone_{opm,n} = Zone_i,
    \]
    \[
    \exists j \in \{n,\dots,Nb_{opm}\} \text{ tel que } 
    \begin{cases}
    Top_{m,n} = \text{Ébauche} \Rightarrow Top_{m,j} \in \{\text{Reprise}, \text{Demi-finition}, \text{Finition}, \text{Super-finition}\} \\
    Top_{m,n} = \text{Reprise} \Rightarrow Top_{m,j} \in \{\text{Demi-finition}, \text{Finition}, \text{Super-finition}\} \\
    Top_{m,n} = \text{Demi-finition} \Rightarrow Top_{m,j} \in \{\text{Finition}, \text{Super-finition}\} \\
    Top_{m,n} = \text{Finition} \Rightarrow Top_{m,j} = \text{Super-finition}
    \end{cases}
    \]
    
    \item \textbf{Contrainte de continuité d'outil (optionnelle)} :
    \[
    \exists n \text{ tel que } Outil_{m,n} = Outil_{m,n+1} \Rightarrow \text{Regroupement possible}
    \]
\end{itemize}

\begin{tcolorbox}[colback=green!5!white, colframe=green!50!black, title=Synthèse : Analyse des Contraintes de Fabrication]
\begin{center}
\textbf{Pour une analyse exhaustive des contraintes, le méthodiste doit identifier :}
\end{center}
\begin{enumerate}
    \item Les \textbf{contraintes géométriques} : accessibilité, rayons, longueurs d'outil, tolérances dimensionnelles et géométriques.
    \item Les \textbf{contraintes technologiques} : matériau, limites outil, limites machine (puissance, couple, vitesses).
    \item Les \textbf{contraintes de séquencement} : antériorités entre zones, hiérarchie des opérations, regroupements possibles.
\end{enumerate}
\textit{Cette analyse structure le processus de décision et garantit que les gammes générées sont à la fois faisables et robustes.}
\end{tcolorbox}
% ============================================
% PARTIE 2 : MÉTHODOLOGIE DE DÉFINITION DES GAMMES PAR CLASSE
% ============================================

\section*{Partie 2 : Méthodologie de Définition des Gammes par Classe}

\subsection*{Introduction}

La méthodologie de définition des gammes doit s'adapter à la nature du produit et au contexte de production. Les trois classes identifiées (A, B, C) répondent à des logiques distinctes : optimisation du flux pour la grande série, modularité pour la moyenne série, flexibilité pour l'unité. Cette partie détaille pour chaque classe les bonnes pratiques et les outils méthodologiques spécifiques.

% ============================================
% CHAPITRE 4 : CLASSE A - PRODUITS STANDARDS, GRANDE SÉRIE
% ============================================

\section{Définition des Gammes pour la Classe A (Produits Standards, Grande Série)}

\begin{tcolorbox}[colback=blue!10!white, colframe=blue!75!black, title=Caractéristiques de la Classe A]
\begin{itemize}
    \item \textbf{Volume} : Grande série, production continue ou par vagues
    \item \textbf{Demande} : Stable, prévisible
    \item \textbf{Produits} : Standards, faible variabilité
    \item \textbf{Objectif principal} : Productivité maximale, coût unitaire minimal
    \item \textbf{Exemples} : Industrie automobile, électroménager, composants standard
\end{itemize}
\end{tcolorbox}

\subsection{Objectif : Maximiser la productivité et la répétabilité}

L'objectif premier des gammes de Classe A est l'optimisation du couple \textbf{productivité / coût}. La répétabilité est la garantie de la performance.

\subsubsection{Principes fondamentaux}

\begin{itemize}
    \item \textbf{Standardisation maximale} : Réduire la diversité des opérations, des outils et des réglages.
    \item \textbf{Élimination des temps morts} : Réduction des temps de changement, des manipulations, des attentes.
    \item \textbf{Stabilité du processus} : Processus capable et maîtrisé (indice $C_{pk} \geq 1.33$).
    \item \textbf{Répétabilité} : Des résultats identiques quelle que soit l'équipe ou l'opérateur.
\end{itemize}

\begin{tcolorbox}[colback=green!5!white, colframe=green!50!black, title=Indicateurs clés pour la Classe A]
\begin{tabular}{|l|c|}
\hline
\textbf{Indicateur} & \textbf{Cible} \\
\hline
Temps de cycle & Minimal \\
Taux d'utilisation machine & $\geq 85\%$ \\
Taux de rendement synthétique (TRS) & $\geq 85\%$ \\
Temps de changement de série (SMED) & $\leq 10$ min \\
Taux de défauts & $\leq 100$ ppm \\
\hline
\end{tabular}
\end{tcolorbox}

\subsection{Optimisation et figement de la Nomenclature}

\subsubsection{Nomenclature figée}

Pour les produits de grande série, la nomenclature doit être \textbf{figée} et \textbf{gérée par révision}. Toute modification doit suivre un processus rigoureux de validation.

\begin{itemize}
    \item \textbf{Nomenclature de production unique} : Une seule version active à un instant donné.
    \item \textbf{Gestion des révisions} : Historique des modifications tracé, date d'application.
    \item \textbf{Impact analyse} : Évaluation des conséquences sur les stocks, les en-cours, les coûts.
\end{itemize}

\subsubsection{Optimisation de la nomenclature}

\begin{description}
    \item[Rationalisation des composants] Réduire le nombre de références (standardisation).
    \item[Éclatement des nomenclatures] Décomposer pour faciliter la planification et l'approvisionnement.
    \item[Nomenclature de planification] Regrouper les articles en familles pour le Plan Industriel et Commercial (PIC).
\end{description}

\begin{tcolorbox}[colback=yellow!10!white, colframe=yellow!50!black, title=Exemple : Optimisation par standardisation]
\textbf{Avant optimisation :}
\begin{itemize}
    \item 1500 références de vis
    \item 200 références de roulements
    \item 80 références de ressorts
\end{itemize}
\vspace{0.5em}
\textbf{Après optimisation :}
\begin{itemize}
    \item 150 références de vis (standardisation dimensionnelle)
    \item 50 références de roulements (familles identifiées)
    \item 20 références de ressorts (modularité)
\end{itemize}
\textit{Gain : réduction de 70\% du nombre de références, simplification des achats et des stocks.}
\end{tcolorbox}

\subsection{Standardisation des Gammes (Gamme unique et réutilisable)}

\subsubsection{Principe de la gamme standard}

Pour la Classe A, on recherche une \textbf{gamme unique} par article, réutilisable et optimisée. La gamme doit être :

\begin{itemize}
    \item \textbf{Séquentielle} : Opérations enchaînées sans embranchement.
    \item \textbf{Équilibrée} : Temps de cycle constant entre postes.
    \item \textbf{Documentée} : Instructions de travail détaillées, modes opératoires.
    \item \textbf{Paramétrée} : Conditions de coupe, temps unitaires, temps de préparation.
\end{itemize}

\subsubsection{Construction de la gamme standard}

\begin{enumerate}
    \item \textbf{Décomposition en phases} : Regroupement par poste de travail.
    \item \textbf{Identification des opérations élémentaires} : Séquence logique.
    \item \textbf{Détermination des temps} : Chronométrage ou temps préétablis (MTM).
    \item \textbf{Optimisation du cycle} : Suppression des temps masqués, équilibrage.
    \item \textbf{Validation} : Lancement pilote, ajustements, figement.
\end{enumerate}

\subsubsection{Gamme de remplacement (alternative)}

Pour garantir la continuité de production, une \textbf{gamme de remplacement} peut être définie. Elle est utilisée en cas :

\begin{itemize}
    \item D'indisponibilité d'une machine critique
    \item De panne d'outillage spécifique
    \item De variation de charge (sous-traitance)
\end{itemize}

\subsection{Optimisation des Flux (Lots, chevauchement, post-consommation)}

\subsubsection{Taille de lot économique}

La détermination de la taille de lot économique est essentielle pour équilibrer les coûts de lancement et les coûts de stockage.

\begin{tcolorbox}[colback=blue!5!white, colframe=blue!75!black, title=Formule de Wilson (lot économique)]
\[
Q^* = \sqrt{\frac{2 \times D \times C_{\text{lancement}}}{C_{\text{stockage}}}}
\]
Avec :
\begin{itemize}
    \item $Q^*$ : Quantité économique par lot
    \item $D$ : Demande annuelle
    \item $C_{\text{lancement}}$ : Coût de lancement (préparation, réglage)
    \item $C_{\text{stockage}}$ : Coût de possession du stock (\% de la valeur)
\end{itemize}
\end{tcolorbox}

\subsubsection{Chevauchement des opérations}

Pour réduire les temps de cycle totaux, on peut chevaucher les opérations lorsqu'elles sont réalisées sur des postes différents.

\begin{tcolorbox}[colback=gray!5!white, colframe=gray!75!black, title=Principe du chevauchement]
Soit $Q$ la taille du lot, $T_{\text{op1}}$ et $T_{\text{op2}}$ les temps unitaires, et $T_{\text{tr}}$ le temps de transfert entre postes.

Temps total sans chevauchement :
\[
T_{\text{total}} = Q \times (T_{\text{op1}} + T_{\text{op2}})
\]

Temps total avec chevauchement (transfert par lots de $q$ pièces) :
\[
T_{\text{total}} = Q \times T_{\text{op1}} + \frac{Q}{q} \times T_{\text{tr}} + q \times T_{\text{op2}}
\]
\end{tcolorbox}

\subsubsection{Post-consommation}

La post-consommation est une méthode de gestion des matières où la sortie de stock est déclenchée par la déclaration d'achèvement de l'opération.

\begin{description}
    \item[Avantages] :
    \begin{itemize}
        \item Simplification des transactions (pas de sortie manuelle)
        \item Réduction des erreurs de saisie
        \item Meilleure traçabilité des consommations réelles
    \end{itemize}
    \item[Conditions] :
    \begin{itemize}
        \item Processus stabilisé (consommations prévisibles)
        \item Stock de sécurité adapté
        \item Contrôle qualité à la réception
    \end{itemize}
\end{description}

\begin{tcolorbox}[colback=yellow!10!white, colframe=yellow!50!black, title=Mise en œuvre de la post-consommation]
\begin{enumerate}
    \item \textbf{Configuration} : Dans la fiche article, activer l'option "Post-consommer matières".
    \item \textbf{Déclaration d'achèvement} : La sortie matière est automatique lors de la déclaration de fin d'opération.
    \item \textbf{Contrôle} : Comparaison entre consommations estimées et réelles pour détection des écarts.
\end{enumerate}
\textit{Remarque : La post-consommation est particulièrement adaptée aux articles de série où la variabilité est faible.}
\end{tcolorbox}

\subsection{Intégration au Programme Directeur (PDP)}

\subsubsection{Rôle du Programme Directeur de Production}

Le Programme Directeur de Production (PDP) est l'outil de pilotage de la production pour les articles de Classe A. Il traduit le Plan Industriel et Commercial (PIC) en un plan de production détaillé.

\begin{tcolorbox}[colback=blue!5!white, colframe=blue!75!black, title=Articulation PIC - PDP - MRP]
\begin{center}
PIC (Familles de produits, horizons longs)\\
$\downarrow$\\
PDP (Articles finis, horizon moyen)\\
$\downarrow$\\
MRP (Composants, horizon court)
\end{center}
\end{tcolorbox}

\subsubsection{Élaboration du PDP pour la Classe A}

\begin{enumerate}
    \item \textbf{Collecte de la demande} :
    \begin{itemize}
        \item Prévisions de vente
        \item Commandes fermes clients
        \item Stock de sécurité
    \end{itemize}
    
    \item \textbf{Calcul du PDP} :
    \[
    PDP_t = D_t + SS_t - S_{t-1} - E_t
    \]
    Où :
    \begin{itemize}
        \item $PDP_t$ : Programme directeur à la période $t$
        \item $D_t$ : Demande (prévisions + commandes)
        \item $SS_t$ : Stock de sécurité
        \item $S_{t-1}$ : Stock disponible en fin de période précédente
        \item $E_t$ : En-cours de fabrication
    \end{itemize}
    
    \item \textbf{Vérification de la capacité} :
    \begin{itemize}
        \item Confrontation PDP $\leftrightarrow$ capacité disponible
        \item Lissage des charges si dépassement
        \item Validation du PDP
    \end{itemize}
    
    \item \textbf{Lancement en fabrication} :
    \begin{itemize}
        \item Transformation du PDP en ordres de fabrication
        \item Planification des besoins matières (MRP)
        \item Lancement des ordres selon les délais d'obtention
    \end{itemize}
\end{enumerate}

\subsubsection{Horizons et bornes de planification}

Pour garantir la stabilité de la production, des horizons de planification sont définis :

\begin{tabular}{|p{0.3\textwidth}|p{0.65\textwidth}|}
\hline
\textbf{Horizon} & \textbf{Description} \\
\hline
\textbf{Horizon figé} & Période pendant laquelle le PDP ne peut être modifié. Ordres lancés, matières réservées. \\
\hline
\textbf{Horizon de réservation} & Période où les modifications sont possibles mais limitées. Besoins matières réservés. \\
\hline
\textbf{Horizon de planification} & Période de prévision où les ordres peuvent être replanifiés librement. \\
\hline
\end{tabular}

\begin{tcolorbox}[colback=green!5!white, colframe=green!50!black, title=Synthèse : Définition des Gammes pour la Classe A]
\textbf{Points clés à retenir :}
\begin{enumerate}
    \item \textbf{Nomenclature figée} et gérée par révision.
    \item \textbf{Gamme unique standardisée} optimisée pour la productivité.
    \item \textbf{Post-consommation} pour simplifier la gestion des matières.
    \item \textbf{PDP structuré} avec horizons figés pour garantir la stabilité.
    \item \textbf{Indicateurs de performance} : TRS, temps de cycle, taux de défauts.
\end{enumerate}
\vspace{0.5em}
\textit{L'objectif final est d'atteindre un processus stable, répétable et performant, où la variabilité est maîtrisée et les temps morts éliminés.}
\end{tcolorbox}

\subsubsection{Exemple de gamme Classe A : Usinage d'un carter moteur}

\begin{tcolorbox}[colback=gray!5!white, colframe=gray!75!black, title=Gamme simplifiée - Carter moteur (grande série)]
\begin{center}
\begin{tabular}{|c|c|c|c|c|}
\hline
\textbf{Phase} & \textbf{Poste} & \textbf{Opération} & \textbf{Temps unitaire} & \textbf{Temps préparation} \\
\hline
10 & Centre d'usinage 1 & Ébauche face A & 2 min & 30 min \\
20 & Centre d'usinage 1 & Ébauche face B & 2 min & - \\
30 & Centre d'usinage 2 & Perçage (6 trous) & 1.5 min & 15 min \\
40 & Centre d'usinage 2 & Taraudage (6 trous) & 1.5 min & - \\
50 & Contrôle & Contrôle dimensionnel & 3 min & 10 min \\
60 & Lavage & Nettoyage & 2 min & 5 min \\
\hline
\multicolumn{3}{|c|}{\textbf{Temps unitaire total}} & \textbf{12 min} & \textbf{60 min} \\
\hline
\end{tabular}
\vspace{0.5em}

\textit{Lot économique : 500 pièces}\\
\textit{Temps de cycle : 12 min/pièce, cadence : 40 pièces/heure}
\end{center}
\end{tcolorbox}
% ============================================
% CHAPITRE 5 : CLASSE B - PRODUITS À VARIANTES, MOYENNE SÉRIE
% ============================================

\section{Définition des Gammes pour la Classe B (Produits à Variantes, Moyenne Série)}

\begin{tcolorbox}[colback=blueclair, colframe=blue!70!black, title=Caractéristiques de la Classe B]
\begin{itemize}
    \item \textbf{Volume} : Moyenne série, production par lots renouvelables
    \item \textbf{Demande} : Variable, avec des pics et des creux
    \item \textbf{Produits} : Familles de produits avec options et variantes
    \item \textbf{Objectif} : Concilier standardisation et personnalisation
    \item \textbf{Exemples} : Industrie automobile (options), outillage, machines spéciales
\end{itemize}
\end{tcolorbox}

\subsection{Objectif : Concilier standardisation et personnalisation}

La Classe B se situe au cœur du paradoxe industriel : offrir une large gamme de produits personnalisés tout en maintenant des coûts de production proches de ceux de la grande série.

\subsubsection{Les défis de la Classe B}

\begin{itemize}
    \item \textbf{Diversité croissante} : Multiplication des références et des combinaisons possibles
    \item \textbf{Complexité de la planification} : Anticipation des besoins pour chaque variante
    \item \textbf{Gestion des stocks} : Équilibre entre disponibilité et niveau de stock
    \item \textbf{Maîtrise des temps de changement} : Réduction des temps de réglage entre variantes
\end{itemize}

\subsubsection{Principes fondamentaux}

\begin{tcolorbox}[colback=vertclair, colframe=green!50!black, title=Principes de la Classe B]
\begin{enumerate}
    \item \textbf{Modularité} : Conception par modules standardisés combinables
    \item \textbf{Late Differentiation} : Report de la personnalisation au plus tard dans le processus
    \item \textbf{Production en flux tiré} : Fabrication déclenchée par la commande client
    \item \textbf{Gestion des options} : Anticipation des combinaisons fréquentes
\end{enumerate}
\end{tcolorbox}

\subsection{Construction d'une Nomenclature Générique et Modulaire}

\subsubsection{Principe de la nomenclature générique}

La nomenclature générique est une structure arborescente qui représente l'ensemble des composants possibles pour une famille de produits, avec des relations conditionnelles entre eux.

\begin{tcolorbox}[colback=grisclair, colframe=gray!70!black, title=Structure d'une nomenclature générique]
\begin{center}
\textbf{Article Générique : Automobile}
\end{center}
\begin{tabular}{ll}
Niveau 1 & \textbf{Automobile (Générique)} \\
& $\quad$ Niveau 2 : Moteur (Générique) \\
& $\qquad$ Niveau 3 : Moteur essence 1.6L \\
& $\qquad$ Niveau 3 : Moteur essence 2.0L \\
& $\qquad$ Niveau 3 : Moteur diesel 1.6L \\
& $\quad$ Niveau 2 : Transmission (Générique) \\
& $\qquad$ Niveau 3 : Boîte manuelle 5 vitesses \\
& $\qquad$ Niveau 3 : Boîte automatique 6 vitesses \\
& $\quad$ Niveau 2 : Peinture (Générique) \\
& $\qquad$ Niveau 3 : Peinture blanche \\
& $\qquad$ Niveau 3 : Peinture noire \\
& $\qquad$ Niveau 3 : Peinture métallisée \\
\end{tabular}
\end{tcolorbox}

\subsubsection{Modularité et standardisation}

\begin{description}
    \item[Module] Ensemble fonctionnel cohérent, interchangeable avec d'autres modules.
    \item[Interface standardisée] Points de connexion entre modules, définis et figés.
    \item[Options] Choix proposés au client dans le cadre de la configuration.
    \item[Variantes] Combinaisons d'options formant un produit fini.
\end{description}

\subsubsection{Règles de composition}

Les règles de composition définissent les combinaisons possibles et impossibles :

\begin{itemize}
    \item \textbf{Inclusion} : Une option peut nécessiter la présence d'une autre option
    \item \textbf{Exclusion} : Certaines options sont incompatibles entre elles
    \item \textbf{Quantité} : Une option peut avoir un nombre variable d'occurrences
\end{itemize}

\begin{tcolorbox}[colback=orangeclair, colframe=orange!70!black, title=Exemple de règles de composition]
\begin{itemize}
    \item Si "Toit ouvrant" alors "Peinture métallisée" (inclusion)
    \item "Moteur diesel" ET "Boîte automatique" incompatibles (exclusion)
    \item "Kit main courante" = 2 unités (quantité fixe)
\end{itemize}
\end{tcolorbox}

\subsection{Définition d'une Gamme Générique avec opérations conditionnelles}

\subsubsection{Principe de la gamme générique}

La gamme générique décrit l'ensemble des opérations possibles pour fabriquer une famille de produits. Certaines opérations sont conditionnées par le choix de certaines options.

\subsubsection{Structure d'une gamme générique}

\begin{table}[h]
\centering
\begin{tabular}{|c|c|c|c|}
\hline
\textbf{Phase} & \textbf{Poste} & \textbf{Opération} & \textbf{Condition} \\
\hline
10 & Carrosserie & Soudage caisse & Toujours \\
20 & Peinture & Application primaire & Toujours \\
30 & Peinture & Peinture de base & Option "Couleur" \\
31 & Peinture & Peinture métallisée & Option "Peinture métallisée" \\
40 & Montage & Montage moteur & Toujours \\
41 & Montage & Montage turbo & Option "Moteur turbo" \\
50 & Montage & Montage toit ouvrant & Option "Toit ouvrant" \\
60 & Contrôle & Test électronique & Toujours \\
61 & Contrôle & Test spécifique & Option "Pack Premium" \\
\hline
\end{tabular}
\end{table}

\subsubsection{Types d'opérations conditionnelles}

\begin{description}
    \item[Opération obligatoire] Exécutée pour toutes les variantes.
    \item[Opération optionnelle] Exécutée seulement si une option spécifique est sélectionnée.
    \item[Opération alternative] Une opération parmi plusieurs est exécutée selon l'option choisie.
    \item[Opération séquentielle] L'opération est exécutée pour chaque occurrence d'un composant (ex : 4 roues).
\end{description}

\subsection{Lien avec le Configurateur de Produits}

\subsubsection{Rôle du configurateur}

Le configurateur de produits est l'outil qui permet de passer d'un article générique à une variante spécifique, en fonction des choix du client.

\begin{tcolorbox}[colback=blueclair, colframe=blue!70!black, title=Processus de configuration]
\begin{center}
Commande client $\rightarrow$ Sélection des options $\rightarrow$ Configuration $\rightarrow$ Variante produit
\end{center}
\vspace{0.5em}
Le configurateur génère :
\begin{itemize}
    \item La nomenclature spécifique
    \item La gamme spécifique (sélection des opérations conditionnelles)
    \item La structure de prix
    \item Les documents techniques personnalisés
\end{itemize}
\end{tcolorbox}

\subsubsection{Paramètres de configuration}

\begin{itemize}
    \item \textbf{Caractéristiques} : Attributs configurables (couleur, puissance, finition)
    \item \textbf{Options} : Valeurs possibles pour chaque caractéristique
    \item \textbf{Contraintes} : Règles de cohérence entre options
    \item \textbf{Relations avec la gamme} : Correspondance option $\rightarrow$ opération
\end{itemize}

\subsubsection{Intégration dans l'ERP}

\begin{description}
    \item[FLB (Final Line Build)] : Système de commande pour les articles d'assemblage configurés.
    \item[Variante de produit] : Instance configurée stockée dans le système.
    \item[Ordre d'assemblage] : Généré à partir de la variante pour la production.
    \item[Post-consommation] : Consommation des composants déclenchée par l'achèvement.
\end{description}

\subsection{Mise en place du Séquencement de Ligne Mixte}

\subsubsection{Principe du flux mixte}

Le séquencement de ligne mixte consiste à organiser la production sur une même ligne d'assemblage de variantes différentes, en optimisant l'ordre de passage pour minimiser les temps de changement.

\begin{tcolorbox}[colback=grisclair, colframe=gray!70!black, title=Objectifs du séquencement mixte]
\begin{itemize}
    \item Minimiser les temps de changement d'outils et d'accessoires
    \item Lisser la charge sur les postes de travail
    \item Respecter les dates de livraison
    \item Gérer les contraintes de capacité
    \item Maintenir un flux régulier
\end{itemize}
\end{tcolorbox}

\subsubsection{Règles de séquencement}

\begin{description}
    \item[Règles d'assortiment] Proportion de chaque variante dans le séquencement.
    \begin{itemize}
        \item Restriction de capacité : nombre maximum d'une variante par fenêtre
        \item Proportionnel : respect d'un ratio entre variantes
    \end{itemize}
    
    \item[Règles de positionnement] Organisation relative des variantes.
    \begin{itemize}
        \item Mise en cluster : regrouper les variantes avec mêmes caractéristiques
        \item Blocage : empêcher certaines successions de variantes
        \item Priorité : ordre de traitement selon des critères (date de livraison, client)
    \end{itemize}
\end{description}

\subsubsection{Exemple de séquencement}

\begin{tcolorbox}[colback=vertclair, colframe=green!50!black, title=Exemple : Séquencement d'une ligne d'assemblage automobile]
\begin{table}[h]
\centering
\begin{tabular}{|c|c|c|c|c|}
\hline
\textbf{Position} & \textbf{Moteur} & \textbf{Couleur} & \textbf{Toit} & \textbf{Temps changement} \\
\hline
1 & Essence 1.6 & Blanc & Ferme & -- \\
2 & Essence 1.6 & Blanc & Ferme & Faible \\
3 & Essence 1.6 & Rouge & Ferme & Moyen \\
4 & Diesel 1.6 & Rouge & Ferme & Faible \\
5 & Diesel 1.6 & Noir & Ouvrant & Élevé \\
6 & Essence 2.0 & Noir & Ouvrant & Faible \\
7 & Essence 2.0 & Noir & Ferme & Faible \\
\hline
\end{tabular}
\end{table}
\textit{Note : Les positions 3 et 4 (changement moteur) nécessitent un réglage, mais la couleur identique réduit le temps de changement de peinture.}
\end{tcolorbox}

\subsubsection{Outils de gestion du séquencement}

\begin{itemize}
    \item \textbf{Simulation et création de séquences} : Génération automatique de la séquence optimale
    \item \textbf{Réorganisation des lignes} : Réassortiment des ordres selon les règles
    \item \textbf{Tableau de bord} : Visualisation de l'utilisation des lignes
    \item \textbf{Indicateurs de performance} : Taux d'utilisation, nombre de changements, temps d'arrêt
\end{itemize}

\subsubsection{Gestion des tampons et des segments}

\begin{description}
    \item[Tampon] Zone de stockage entre segments de ligne permettant de réordonnancer.
    \item[Segment de ligne] Ensemble de postes consécutifs entre deux tampons.
    \item[Emplacement de mémoire vive] Capacité du tampon à stocker des produits en attente.
\end{description}

\subsection{Indicateurs de performance pour la Classe B}

\begin{table}[h]
\centering
\begin{tabular}{|l|c|}
\hline
\textbf{Indicateur} & \textbf{Cible} \\
\hline
Taux de personnalisation & Variable selon marché \\
Délai de configuration & 1 à 5 minutes \\
Temps de changement entre variantes & 1 à 5 minutes \\
Taux de service & 98\% \\
Niveau de modularité & 70\% de composants communs \\
\hline
\end{tabular}
\end{table}

\begin{tcolorbox}[colback=blueclair, colframe=blue!70!black, title=Synthèse : Classe B]
\begin{enumerate}
    \item \textbf{Nomenclature générique} : Structure modulaire avec règles de composition
    \item \textbf{Gamme générique} : Opérations conditionnelles liées aux options
    \item \textbf{Configurateur} : Interface entre commande client et production
    \item \textbf{Séquencement mixte} : Organisation optimale de la production multi-variantes
    \item \textbf{Objectif} : Mass Customization - production de masse personnalisée
\end{enumerate}
\vspace{0.5em}
\textit{L'enjeu de la Classe B est de proposer une large variété de produits tout en maîtrisant la complexité industrielle et les coûts de production.}
\end{tcolorbox}
% ============================================
% CHAPITRE 6 : CLASSE C - PRODUITS SPÉCIFIQUES, PETITE SÉRIE / UNITAIRE
% ============================================

\section{Définition des Gammes pour la Classe C (Produits Spécifiques, Petite Série / Unitaire)}

\begin{tcolorbox}[colback=blueclair, colframe=blue!70!black, title=Caractéristiques de la Classe C]
\begin{itemize}
    \item \textbf{Volume} : Petite série, production unitaire
    \item \textbf{Demande} : Unique, non récurrente
    \item \textbf{Produits} : Spécifiques, conçus pour un client ou un projet
    \item \textbf{Objectif} : Assurer la faisabilité et maîtriser les risques
    \item \textbf{Exemples} : Aéronautique (pièces de structure), outillage lourd, machines spéciales, prototypes
\end{itemize}
\end{tcolorbox}

\subsection{Objectif : Assurer la faisabilité et contrôler les risques projet}

La Classe C se caractérise par une production unitaire ou en très petite série, où chaque pièce est unique et souvent complexe. L'objectif n'est plus l'optimisation marginale du temps de cycle, mais la garantie de faisabilité et la maîtrise des risques.

\subsubsection{Les spécificités de la Classe C}

\begin{itemize}
    \item \textbf{Produit unique} : Chaque commande est un projet avec ses propres spécifications
    \item \textbf{Matériaux complexes} : Alliages réfractaires (titane, inconel), composites
    \item \textbf{Géométries complexes} : Formes aérodynamiques, tolérances serrées
    \item \textbf{Coût unitaire élevé} : Valeur ajoutée importante, réduction des risques de rebut
    \item \textbf{Délais contraints} : Calendrier de livraison souvent fixé contractuellement
    \item \textbf{Traçabilité renforcée} : Exigences qualité et certification
\end{itemize}

\subsubsection{Principes fondamentaux}

\begin{tcolorbox}[colback=vertclair, colframe=green!50!black, title=Principes de la Classe C]
\begin{enumerate}
    \item \textbf{Approche projet} : Organisation par projet avec chef de projet dédié
    \item \textbf{Ingénierie concourante} : Collaboration étroite conception / méthodes / production
    \item \textbf{Gestion des risques} : Identification, quantification et mitigation des risques
    \item \textbf{Traçabilité totale} : Suivi complet de la matière, des opérations et des contrôles
    \item \textbf{Flexibilité} : Capacité d'adaptation face aux aléas et modifications
\end{enumerate}
\end{tcolorbox}

\subsection{Création d'une structure par Projet (PCS) pour chaque commande}

\subsubsection{Principe du Projet Client Spécifique (PCS)}

Le PCS est une structure de gestion qui permet de regrouper l'ensemble des données et des activités liées à une commande client spécifique.

\begin{tcolorbox}[colback=grisclair, colframe=gray!70!black, title=Structure d'un projet PCS]
\begin{center}
\textbf{Projet Principal}
\end{center}
\begin{tabular}{ll}
& $\quad$ \textbf{Sous-projet 1 : Conception} \\
& $\qquad$ Activité 1.1 : Étude de définition \\
& $\qquad$ Activité 1.2 : Validation CAO \\
& $\quad$ \textbf{Sous-projet 2 : Méthodes} \\
& $\qquad$ Activité 2.1 : Définition des gammes \\
& $\qquad$ Activité 2.2 : Choix des outils \\
& $\quad$ \textbf{Sous-projet 3 : Fabrication} \\
& $\qquad$ Activité 3.1 : Approvisionnement matières \\
& $\qquad$ Activité 3.2 : Usinage ébauche \\
& $\qquad$ Activité 3.3 : Contrôle intermédiaire \\
& $\qquad$ Activité 3.4 : Usinage finition \\
& $\quad$ \textbf{Sous-projet 4 : Contrôle} \\
& $\qquad$ Activité 4.1 : Contrôle dimensionnel \\
& $\qquad$ Activité 4.2 : Contrôle métallurgique \\
& $\qquad$ Activité 4.3 : Établissement certificat \\
\end{tabular}
\end{tcolorbox}

\subsubsection{Création de la structure projet}

\begin{enumerate}
    \item \textbf{Réception de la commande client} : Spécifications techniques, quantités, délais
    \item \textbf{Génération de la structure projet} : Création du projet principal et des sous-projets
    \item \textbf{Rattachement des articles} : Les articles spécifiques sont liés au projet
    \item \textbf{Définition du phasage} : Activités, jalons, relations entre activités
    \item \textbf{Planification réseau} : Calcul des dates au plus tôt et au plus tard
\end{enumerate}

\subsubsection{Types de projets PCS}

\begin{table}[h]
\centering
\begin{tabular}{|p{0.35\textwidth}|p{0.6\textwidth}|}
\hline
\textbf{Type} & \textbf{Description} \\
\hline
Projet principal & Regroupe plusieurs sous-projets, pilotage global \\
Sous-projet & Fait partie d'un projet principal, activités détaillées \\
Projet unique & Projet indépendant non subdivisé \\
Budget & Utilisé pour la planification sans exécution \\
\hline
\end{tabular}
\end{table}

\subsubsection{Articles spécifiques et articles standard à la commande}

\begin{description}
    \item[Article spécifique] Créé spécifiquement pour un projet, avec sa propre nomenclature et gamme.
    \item[Article standard à la commande (STO)] Article standard dont la fabrication est déclenchée par commande, sans création de nouvelle référence.
\end{description}

\subsection{Élaboration d'une Gamme de Fabrication sur Mesure (spécifique)}

\subsubsection{Principe de la gamme spécifique}

La gamme spécifique est élaborée pour un projet particulier, en adaptant des gammes standards existantes ou en créant de toutes pièces des opérations.

\subsubsection{Démarche d'élaboration}

\begin{enumerate}
    \item \textbf{Analyse géométrique} : Identification des zones à usiner, des surfaces fonctionnelles
    \item \textbf{Définition des posages} : Choix des orientations et des montages
    \item \textbf{Sélection des outils} : Choix des outils de coupe adaptés aux matériaux et géométries
    \item \textbf{Détermination des conditions de coupe} : Calcul ou adaptation selon retour d'expérience
    \item \textbf{Définition du séquencement} : Ordre des opérations respectant les antériorités
    \item \textbf{Identification des contrôles} : Points de contrôle dimensionnel et métallurgique
    \item \textbf{Évaluation des risques} : Identification des opérations critiques
\end{enumerate}

\subsubsection{Structure d'une gamme spécifique}

\begin{tcolorbox}[colback=orangeclair, colframe=orange!70!black, title=Exemple : Gamme spécifique pour pièce aéronautique]
\begin{tabular}{|c|c|c|c|c|c|}
\hline
\textbf{Phase} & \textbf{Opération} & \textbf{Poste} & \textbf{Outil} & \textbf{Temps estimé} & \textbf{Risque} \\
\hline
10 & Contrôle matière & Contrôle & -- & 2 h & Faible \\
20 & Tracage & Méthodes & -- & 4 h & Faible \\
30 & Ébauche face A & Fraiseuse & Fraise ø80 & 12 h & Moyen \\
40 & Ébauche face B & Fraiseuse & Fraise ø80 & 10 h & Moyen \\
50 & Contrôle intermédiaire & Contrôle & MMT & 3 h & Faible \\
60 & Demi-finition & Fraiseuse & Fraise ø50 & 8 h & Élevé \\
70 & Finition & Fraiseuse & Fraise ø32 & 6 h & Élevé \\
80 & Perçage & Perceuse & Foret ø10 & 2 h & Faible \\
90 & Contrôle final & Contrôle & MMT & 4 h & Faible \\
100 & Livraison & Expédition & -- & 1 h & Faible \\
\hline
\end{tabular}
\end{tcolorbox}

\subsection{Intégration de la modélisation et de la quantification des Risques}

\subsubsection{Identification des risques en usinage}

\begin{table}[h]
\centering
\begin{tabular}{|p{0.35\textwidth}|p{0.6\textwidth}|}
\hline
\textbf{Type de risque} & \textbf{Cause possible} \\
\hline
Rupture d'outil & Conditions de coupe excessives, hétérogénéité matière \\
Usure excessive & Durée de vie dépassée, lubrification insuffisante \\
Vibrations & Engagement radial mal maîtrisé, rigidité insuffisante \\
Défaut géométrique & Flexion de l'outil, mauvais positionnement \\
Non-conformité matière & Défaut de forge, inclusion, traitement inadapté \\
Délai dépassé & Aléas de production, rupture d'outil, panne machine \\
\hline
\end{tabular}
\end{table}

\subsubsection{Modélisation du risque outil}

La flexion de l'outil est un indicateur de risque majeur en usinage de précision.

\begin{tcolorbox}[colback=blueclair, colframe=blue!70!black, title=Calcul de la flexion de l'outil]
\[
f_{n,m} = \frac{1280}{\pi^2} \times \frac{K_c}{E} \times F_{z_{n,m}} \times Z_{n,m} \times Ap_{n,m} \times Ae_{n,m} \times \frac{L_{n,m}^3}{D_{n,m}^3}
\]
Avec :
\begin{itemize}
    \item $f$ : Flèche de l'outil (mm)
    \item $K_c$ : Force de coupe spécifique du matériau
    \item $E$ : Module d'Young de l'outil
    \item $F_z$ : Avance à la dent (mm/dent)
    \item $Z$ : Nombre de dents
    \item $Ap$ : Profondeur de passe axiale (mm)
    \item $Ae$ : Engagement radial (mm)
    \item $L$ : Longueur de l'outil (mm)
    \item $D$ : Diamètre de l'outil (mm)
\end{itemize}
\end{tcolorbox}

\subsubsection{Indice de risque par opération}

\[
IR_{\text{op}} = \frac{F_z \times Z \times Ap \times Ae \times \frac{L^3}{D^3}}{R_{\text{seuil}}}
\]

\subsubsection{Attitudes face au risque}

La perception du risque varie selon le décideur. Quatre attitudes peuvent être modélisées :

\begin{table}[h]
\centering
\begin{tabular}{|p{0.25\textwidth}|p{0.7\textwidth}|}
\hline
\textbf{Attitude} & \textbf{Formulation} \\
\hline
Optimiste & $\displaystyle IR_{\text{gamme}} = \min_n(IR_{\text{op}})$ \\
Neutre & $\displaystyle IR_{\text{gamme}} = \frac{1}{N} \sum_{n=1}^N IR_{\text{op}}$ \\
Pessimiste & $\displaystyle IR_{\text{gamme}} = \max_n(IR_{\text{op}})$ \\
Robuste & $\displaystyle S/N = -10 \log\left(S^2 + \frac{1}{\overline{Y}^2}\right)$ \\
\hline
\end{tabular}
\end{tcolorbox}

\subsubsection{Grille de criticité des risques}

\begin{tcolorbox}[colback=grisclair, colframe=gray!70!black, title=Matrice de criticité]
\begin{center}
\begin{tabular}{|c|c|c|c|c|c|}
\hline
\multicolumn{2}{|c|}{} & \multicolumn{4}{|c|}{\textbf{Gravité}} \\
\cline{3-6}
\multicolumn{2}{|c|}{} & \textbf{Mineure} & \textbf{Significative} & \textbf{Critique} & \textbf{Catastrophique} \\
\hline
\multirow{4}{*}{\rotatebox{90}{\textbf{Probabilité}}} & Très faible & Faible & Faible & Modéré & Élevé \\
\cline{2-6}
& Faible & Faible & Modéré & Élevé & Élevé \\
\cline{2-6}
& Modérée & Modéré & Élevé & Critique & Critique \\
\cline{2-6}
& Élevée & Élevé & Critique & Critique & Critique \\
\hline
\end{tabular}
\end{center}
\textit{Seuils d'acceptabilité : Faible = acceptable, Modéré = sous surveillance, Élevé = plan de mitigation, Critique = non acceptable}
\end{tcolorbox}

\subsection{Planification par Contraintes et gestion des goulots}

\subsubsection{Principe de la planification par contraintes}

La planification par contraintes (ou \textit{Theory of Constraints}) consiste à identifier le goulet d'étranglement du processus et à organiser la production autour de cette contrainte.

\begin{tcolorbox}[colback=vertclair, colframe=green!50!black, title=La méthode des 5 étapes (Goldratt)]
\begin{enumerate}
    \item \textbf{Identifier} la contrainte du système
    \item \textbf{Exploiter} la contrainte (maximiser son utilisation)
    \item \textbf{Subordonner} toutes les autres ressources à la contrainte
    \item \textbf{Élever} la contrainte (augmenter sa capacité)
    \item \textbf{Répéter} le processus (la contrainte s'est déplacée)
\end{enumerate}
\end{tcolorbox}

\subsubsection{Identification des goulots pour la Classe C}

Dans un contexte de production unitaire, les goulots sont souvent :

\begin{itemize}
    \item Les machines spécifiques (fraiseuse 5 axes, centre d'usinage grande dimension)
    \item Les opérateurs hautement qualifiés (réglage, programmation)
    \item Les outils de mesure (MMT, contrôles non destructifs)
    \item Les approvisionnements matières (délais longs)
\end{itemize}

\subsubsection{Gestion des contraintes de capacité}

\begin{table}[h]
\centering
\begin{tabular}{|l|l|}
\hline
\textbf{Type de contrainte} & \textbf{Action} \\
\hline
Contrainte machine & Planifier en priorité, amortir les aléas, préparer les outils \\
Contrainte opérateur & Anticiper les disponibilités, former des polyvalents \\
Contrainte outil & Gérer le stock, prévoir les outils frères \\
Contrainte matière & Anticiper les commandes, gérer les stocks de sécurité \\
\hline
\end{tabular}
\end{table}

\subsubsection{Planification réseau de projet}

La planification réseau permet de modéliser les dépendances entre activités et de calculer les marges.

\begin{tcolorbox}[colback=blueclair, colframe=blue!70!black, title=Exemple de planification réseau]
\begin{center}
\begin{tabular}{lccccc}
\hline
\textbf{Activité} & \textbf{Durée (j)} & \textbf{Début au plus tôt} & \textbf{Fin au plus tôt} & \textbf{Début au plus tard} & \textbf{Marge} \\
\hline
A : Conception & 5 & 0 & 5 & 0 & 0 \\
B : Méthodes & 3 & 5 & 8 & 7 & 2 \\
C : Appro. matières & 10 & 5 & 15 & 5 & 0 \\
D : Usinage & 8 & 15 & 23 & 15 & 0 \\
E : Contrôle & 2 & 23 & 25 & 23 & 0 \\
\hline
\end{tabular}
\end{center}
\textit{Chemin critique : A - C - D - E (durée totale 25 jours)}
\end{tcolorbox}

\subsubsection{Indicateurs de pilotage pour la Classe C}

\begin{table}[h]
\centering
\begin{tabular}{|l|c|}
\hline
\textbf{Indicateur} & \textbf{Objectif} \\
\hline
Respect du délai projet & 100\% \\
Taux de première pièce conforme & $\geq$ 95\% \\
Taux de rebut & $\leq$ 2\% \\
Respect du budget & $\pm$ 5\% \\
Taux de traçabilité & 100\% \\
\hline
\end{tabular}
\end{table}

\subsubsection{Gestion des modifications et des aléas}

\begin{itemize}
    \item \textbf{Suivi des modifications} : Gestion des évolutions de conception en cours de projet
    \item \textbf{Plan de continuité} : Alternatives en cas de panne ou d'indisponibilité
    \item \textbf{Revues de projet} : Points d'avancement réguliers avec les parties prenantes
    \item \textbf{Retour d'expérience} : Capitalisation des enseignements pour les projets futurs
\end{itemize}

\begin{tcolorbox}[colback=blueclair, colframe=blue!70!black, title=Synthèse : Classe C]
\begin{enumerate}
    \item \textbf{Structure projet (PCS)} : Organisation par projet avec activités et jalons
    \item \textbf{Gamme spécifique} : Adaptée à chaque pièce, avec opérations détaillées
    \item \textbf{Gestion des risques} : Identification, modélisation et quantification des risques
    \item \textbf{Planification par contraintes} : Pilotage par les goulets et le chemin critique
    \item \textbf{Objectif} : Fiabiliser la production unitaire et maîtriser les risques projet
\end{enumerate}
\vspace{0.5em}
\textit{L'enjeu de la Classe C est de garantir la faisabilité et la qualité d'une production unique, avec une maîtrise totale des risques et des délais.}
\end{tcolorbox}
% ============================================
% PARTIE 3 : FORMALISATION, OPTIMISATION ET INTÉGRATION ERP
% ============================================

\section*{Partie 3 : Formalisation, Optimisation et Intégration ERP}

\subsection*{Introduction}

La formalisation des données de production est une étape critique pour assurer l'intégration cohérente entre les différents modules de l'ERP. Cette partie détaille la structuration des informations techniques et leur saisie dans le système, garantissant ainsi la fiabilité et la réutilisabilité des données.

% ============================================
% CHAPITRE 7 : FORMALISATION DES DONNÉES POUR L'ERP
% ============================================

\section{Formalisation des Données pour l'ERP}

\subsection{Structuration des informations (Articles, Nomenclatures, Gammes, Ressources)}

La structuration des données dans l'ERP repose sur une organisation hiérarchique et cohérente des informations techniques.

\subsubsection{Données Articles}

\begin{tcolorbox}[colback=blueclair, colframe=blue!70!black, title=Structure type d'une fiche article ERP]
\begin{tabular}{|l|l|}
\hline
\textbf{Champ} & \textbf{Description} \\
\hline
Code article & Identifiant unique (ex: \textsc{FLT-1000-A}) \\
Désignation & Libellé descriptif (ex: Fraise torique diamètre 10) \\
Type d'article & Acheté / Fabriqué / Générique / Produit / Outil \\
Système de commande & SIC / Planifié / FLB / Manuel \\
Unité de stock & Pièce / Kg / Mètre \\
Poids & Poids unitaire \\
Géré par lot & Oui / Non \\
Sérialisé & Oui / Non \\
Méthode de valorisation & Coût standard / Prix de la série / FIFO \\
\hline
\end{tabular}
\end{tcolorbox}

\subsubsection{Données Nomenclatures}

\begin{table}[h]
\centering
\begin{tabular}{|p{0.3\textwidth}|p{0.65\textwidth}|}
\hline
\textbf{Champ} & \textbf{Description} \\
\hline
Article parent & Code de l'article assemblé \\
Article composant & Code de l'article constitutif \\
Quantité & Nombre d'unités du composant \\
Coef. de rebut & Pourcentage de perte prévisionnelle \\
Date d'application & Date de début de validité \\
Date d'expiration & Date de fin de validité \\
Opération de rattachement & Liens avec la gamme (phase d'utilisation) \\
\hline
\end{tabular}
\end{table}

\begin{tcolorbox}[colback=grisclair, colframe=gray!70!black, title=Exemple : Nomenclature de production - Brouette]
\begin{tabular}{|l|l|c|}
\hline
\textbf{Article parent} & \textbf{Composant} & \textbf{Quantité} \\
\hline
\multirow{5}{*}{BROUETTE-001} & SUPPORT-001 & 1 \\
 & BENNE-001 & 1 \\
 & ROU-001 & 1 \\
 & AXE-001 & 1 \\
 & PEINTURE-001 & 0,5 m² \\
\hline
\end{tabular}
\end{tcolorbox}

\subsubsection{Données Gammes}

\begin{table}[h]
\centering
\begin{tabular}{|p{0.25\textwidth}|p{0.7\textwidth}|}
\hline
\textbf{Champ} & \textbf{Description} \\
\hline
Code gamme & Identifiant de la gamme \\
Article fabriqué & Article auquel la gamme est associée \\
Site & Site de production \\
Phase & Numéro d'ordre (10, 20, 30...) \\
Poste de charge & Centre de travail utilisé \\
Opération & Description de l'opération \\
Temps de préparation & Temps de réglage (hors série) \\
Temps unitaire & Temps par pièce \\
Outil & Référence de l'outil à utiliser \\
\hline
\end{tabular}
\end{table}

\begin{tcolorbox}[colback=vertclair, colframe=green!50!black, title=Exemple : Gamme d'usinage - Fraisage]
\begin{tabular}{|c|c|c|c|c|c|}
\hline
\textbf{Phase} & \textbf{Poste} & \textbf{Opération} & \textbf{Tps prép.} & \textbf{Tps unitaire} & \textbf{Outil} \\
\hline
10 & FRAISE-3AX & Ébauche face A & 30 min & 12 min & FRAISE-080 \\
20 & FRAISE-3AX & Ébauche face B & -- & 10 min & FRAISE-080 \\
30 & FRAISE-3AX & Finition & 15 min & 8 min & FRAISE-032 \\
40 & CONTROLE & Contrôle final & 10 min & 5 min & PALPEUR \\
\hline
\end{tabular}
\end{tcolorbox}

\subsubsection{Données Ressources}

\begin{table}[h]
\centering
\begin{tabular}{|p{0.3\textwidth}|p{0.65\textwidth}|}
\hline
\textbf{Champ} & \textbf{Description} \\
\hline
Code ressource & Identifiant du poste de charge \\
Désignation & Libellé descriptif \\
Type & Machine / Atelier / Poste manuel \\
Capacité & Nombre d'unités simultanées \\
Calendrier & Code calendrier associé \\
Taux horaire & Coût horaire machine \\
\hline
\end{tabular}
\end{table}

\subsection{Gestion des révisions et des modifications techniques (ECN)}

\subsubsection{Principe de la gestion des révisions}

Les modifications techniques doivent être tracées et contrôlées pour garantir la cohérence des données et la traçabilité des évolutions.

\begin{tcolorbox}[colback=blueclair, colframe=blue!70!black, title=Cycle de vie d'une révision]
\begin{center}
Créé $\rightarrow$ En validation $\rightarrow$ Approuvé $\rightarrow$ Actif $\rightarrow$ Remplacé $\rightarrow$ Archivé
\end{center}
\end{tcolorbox}

\subsubsection{Structure d'une révision}

\begin{table}[h]
\centering
\begin{tabular}{|l|l|}
\hline
\textbf{Champ} & \textbf{Description} \\
\hline
Numéro de révision & 001, 002, 003... \\
Statut & Nouveau / Approuvé / Actif / Remplacé \\
Date d'application & Date de début de validité \\
Date d'expiration & Date de fin de validité \\
Motif & Description de la modification \\
Auteur & Personne ayant créé la révision \\
Approbateur & Personne ayant validé la révision \\
\hline
\end{tabular}
\end{table}

\subsubsection{Processus de modification technique (ECN)}

\begin{enumerate}
    \item \textbf{Demande de modification} : Identification du besoin (non-conformité, amélioration)
    \item \textbf{Analyse d'impact} : Évaluation des conséquences (stocks, en-cours, coûts)
    \item \textbf{Création de la révision} : Saisie des nouvelles données dans l'ERP
    \item \textbf{Validation} : Approbation par les responsables concernés
    \item \textbf{Activation} : Mise en application à la date prévue
    \item \textbf{Gestion des stocks} : Traitement des stocks résiduels de l'ancienne version
\end{enumerate}

\begin{tcolorbox}[colback=orangeclair, colframe=orange!70!black, title=Exemple : Modification technique d'un outil]
\begin{tabular}{|l|l|}
\hline
\textbf{Élément} & \textbf{Information} \\
\hline
Révision & 002 \\
Article & FRAISE-080 \\
Modification & Changement du revêtement \\
Motif & Augmentation de la durée de vie (12h $\rightarrow$ 18h) \\
Date d'application & 01/06/2024 \\
Gestion des stocks & Écoulement de la révision 001 jusqu'au 31/05/2024 \\
\hline
\end{tabular}
\end{tcolorbox}

\subsubsection{Gestion des modifications dans la nomenclature}

\begin{itemize}
    \item \textbf{Remplacement} : Un composant est remplacé par un autre
    \item \textbf{Ajout} : Un nouveau composant est introduit
    \item \textbf{Suppression} : Un composant est retiré
    \item \textbf{Modification de quantité} : La quantité d'un composant change
    \item \textbf{Date d'application} : Date à partir de laquelle la modification s'applique
\end{itemize}

\subsection{Saisie des données dans les fiches ERP (exemples concrets)}

\subsubsection{Fiche Article : Fraise torique diamètre 10 mm}

\begin{tcolorbox}[colback=grisclair, colframe=gray!70!black, title=Exemple : Fiche article ERP - Fraise torique]
\begin{tabular}{|l|l|}
\hline
\textbf{Champ} & \textbf{Valeur} \\
\hline
Code article & \textsc{FR-TOR-010} \\
Désignation & Fraise torique diamètre 10, rayon 2 \\
Type d'article & Acheté \\
Système de commande & SIC \\
Unité de stock & Pièce \\
Poids & 0,05 kg \\
Géré par lot & Non \\
Sérialisé & Non \\
Méthode de valorisation & Coût standard \\
Coût standard & 85,00 € \\
Fournisseur privilégié & SANDVIK \\
Délai d'approvisionnement & 5 jours \\
Stock de sécurité & 10 \\
Point de commande & 25 \\
\hline
\end{tabular}
\end{tcolorbox}

\subsubsection{Fiche Article : Carter moteur fabriqué}

\begin{tcolorbox}[colback=vertclair, colframe=green!50!black, title=Exemple : Fiche article ERP - Carter moteur]
\begin{tabular}{|l|l|}
\hline
\textbf{Champ} & \textbf{Valeur} \\
\hline
Code article & \textsc{CARTER-M12} \\
Désignation & Carter moteur M12 \\
Type d'article & Fabriqué \\
Système de commande & Planifié \\
Unité de stock & Pièce \\
Poids & 12,5 kg \\
Géré par lot & Oui \\
Sérialisé & Oui \\
Méthode de valorisation & Coût standard \\
Gamme par défaut & \textsc{GAM-CARTER-M12} \\
Délai de fabrication & 3 jours \\
\hline
\end{tabular}
\end{tcolorbox}

\subsubsection{Fiche Nomenclature - Niveau 1}

\begin{tcolorbox}[colback=blueclair, colframe=blue!70!black, title=Exemple : Nomenclature - Carter moteur M12]
\begin{tabular}{|l|l|c|c|}
\hline
\textbf{Composant} & \textbf{Désignation} & \textbf{Quantité} & \textbf{Opération} \\
\hline
\textsc{BRUT-CARTER} & Ébauche carter & 1 & 10 \\
\textsc{JOINT-012} & Joint d'étanchéité & 1 & 50 \\
\textsc{ECROU-M8} & Écrou M8 & 6 & 60 \\
\textsc{PEINTURE} & Peinture haute température & 0,2 & 70 \\
\hline
\end{tabular}
\end{tcolorbox}

\subsubsection{Fiche Gamme - Carter moteur M12}

\begin{tcolorbox}[colback=grisclair, colframe=gray!70!black, title=Exemple : Gamme de fabrication - Carter moteur M12]
\begin{tabular}{|c|c|c|c|c|}
\hline
\textbf{Phase} & \textbf{Poste} & \textbf{Opération} & \textbf{Tps prép.} & \textbf{Tps unitaire} \\
\hline
10 & FRAISE-5AX & Ébauche carter & 45 min & 25 min \\
20 & FRAISE-5AX & Perçage (6 trous) & 15 min & 8 min \\
30 & FRAISE-5AX & Taraudage M8 & 10 min & 6 min \\
40 & CONTROLE & Contrôle dimensionnel & 20 min & 10 min \\
50 & LAVAGE & Nettoyage & 10 min & 5 min \\
60 & PEINTURE & Peinture & 15 min & 8 min \\
70 & CONTRÔLE & Contrôle final & 10 min & 4 min \\
\hline
\end{tabular}
\end{tcolorbox}

\subsubsection{Fiche Poste de Charge}

\begin{tcolorbox}[colback=orangeclair, colframe=orange!70!black, title=Exemple : Fiche poste de charge - FRAISE-5AX]
\begin{tabular}{|l|l|}
\hline
\textbf{Champ} & \textbf{Valeur} \\
\hline
Code poste & \textsc{FRAISE-5AX} \\
Désignation & Centre d'usinage 5 axes \\
Type & Machine \\
Capacité & 1 machine \\
Calendrier & 3x8 (lundi au vendredi) \\
Taux horaire machine & 95,00 €/h \\
Taux horaire main-d'œuvre & 45,00 €/h \\
Puissance broche & 100 kW \\
Vitesse max & 1500 tr/min \\
Magasin outil & 60 positions \\
\hline
\end{tabular}
\end{tcolorbox}

\subsubsection{Fiche Outil}

\begin{tcolorbox}[colback=vertclair, colframe=green!50!black, title=Exemple : Fiche outil - Fraise torique]
\begin{tabular}{|l|l|}
\hline
\textbf{Champ} & \textbf{Valeur} \\
\hline
Code outil & \textsc{FR-TOR-010} \\
Désignation & Fraise torique Ø10 R2 \\
Type outil & Fraise torique \\
Diamètre & 10 mm \\
Rayon de bec & 2 mm \\
Longueur & 80 mm \\
Nombre de dents & 4 \\
Nombre de plaquettes & 4 \\
Coût par arête & 12,50 € \\
Durée de vie & 120 min \\
Vitesse de coupe & 120 m/min \\
Avance à la dent & 0,12 mm/dent \\
Profondeur max & 5 mm \\
Engagement radial max & 8 mm \\
\hline
\end{tabular}
\end{tcolorbox}

\subsubsection{Saisie des données : Bonnes pratiques}

\begin{tcolorbox}[colback=blueclair, colframe=blue!70!black, title=Recommandations pour la saisie ERP]
\begin{itemize}
    \item \textbf{Standardisation des codifications} : Utiliser des masques de saisie cohérents
    \item \textbf{Validation des données} : Vérifier la cohérence des données avant validation
    \item \textbf{Traçabilité} : Renseigner systématiquement l'auteur et la date de création
    \item \textbf{Révision} : Gérer les évolutions par révision, pas par écrasement
    \item \textbf{Documentation} : Joindre les documents techniques (plans, instructions)
    \item \textbf{Sécurité} : Gérer les droits d'accès par profil (méthodes, production, achats)
\end{itemize}
\end{tcolorbox}

\begin{tcolorbox}[colback=gray!5!white, colframe=gray!75!black, title=Synthèse : Formalisation des Données pour l'ERP]
\begin{enumerate}
    \item \textbf{Articles} : Fiche complète avec type, système de commande, données techniques
    \item \textbf{Nomenclatures} : Structure arborescente avec quantités et opérations de rattachement
    \item \textbf{Gammes} : Séquence des opérations avec temps de préparation et temps unitaires
    \item \textbf{Ressources} : Postes de charge avec capacités, calendriers et taux horaires
    \item \textbf{Révisions} : Gestion des modifications techniques avec cycle de validation
\end{enumerate}
\vspace{0.5em}
\textit{La qualité des données saisies dans l'ERP conditionne la fiabilité de la planification, le calcul des coûts et la traçabilité de la production.}
\end{tcolorbox}
% ============================================
% CHAPITRE 8 : OPTIMISATION DES GAMMES
% ============================================

\section{Optimisation des Gammes}

\subsection{Définition des Indicateurs de Performance (Temps VA/NVA, coût outil, efficacité)}

L'optimisation des gammes nécessite une définition précise des indicateurs de performance permettant de comparer objectivement différentes solutions.

\subsubsection{Temps à Valeur Ajoutée (VA)}

Le temps à Valeur Ajoutée correspond au temps pendant lequel la machine ou l'opérateur génère de la valeur sur la pièce.

\begin{tcolorbox}[colback=blueclair, colframe=blue!70!black, title=Calcul du temps VA]
\[
T_{VA} = \sum_{n=1}^{N} T_{\text{usinage},n}
\]
Avec :
\begin{itemize}
    \item $T_{\text{usinage},n} = \dfrac{L_{\text{trajet},n} \times N_{\text{passe axiale}} \times N_{\text{passe radiale}}}{V_{f,n}}$
    \item $L_{\text{trajet}}$ : Longueur du trajet élémentaire
    \item $V_f$ : Vitesse d'avance de l'outil
\end{itemize}
\end{tcolorbox}

\subsubsection{Temps à Non-Valeur Ajoutée (NVA)}

Les temps NVA regroupent tous les temps pendant lesquels l'outil ne coupe pas.

\begin{table}[h]
\centering
\begin{tabular}{|p{0.4\textwidth}|p{0.55\textwidth}|}
\hline
\textbf{Type de temps NVA} & \textbf{Formule} \\
\hline
Changement d'orientation & $\displaystyle T_{\text{orientation}} = \sum_{n=2}^{N} V_{\alpha_n} \times T_{\text{unit orientation}}$ \\
Changement d'accessoire & $\displaystyle T_{\text{accessoire}} = \sum_{n=2}^{N} X_{\text{acc},n} \times T_{\text{unit acc}}$ \\
Changement d'outil & $\displaystyle T_{\text{outil}} = \sum_{n=2}^{N} Y_{\text{outil},n} \times T_{\text{unit outil}}$ \\
Changement de plaquette & $\displaystyle T_{\text{plaquette}} = \sum_{n=1}^{N} \left\lfloor \frac{T_{\text{usinage},n}}{DDV_{\text{outil}}} \right\rfloor \times N_{\text{plaquettes}} \times T_{\text{unit plaq}}$ \\
\hline
\end{tabular}
\end{table}

\subsubsection{Coût outil (coût carbure)}

Le coût outil représente la part du coût de fabrication liée à l'usure et au remplacement des outils coupants.

\begin{tcolorbox}[colback=grisclair, colframe=gray!70!black, title=Calcul du coût outil]
\[
C_{\text{outil}} = \sum_{n=1}^{N} \left\lfloor \frac{T_{\text{usinage},n}}{DDV_{\text{outil},n}} \right\rfloor \times N_{\text{plaquettes},n} \times C_{\text{arête},n}
\]
Avec :
\begin{itemize}
    \item $DDV$ : Durée de vie de l'outil (min)
    \item $N_{\text{plaquettes}}$ : Nombre de plaquettes de l'outil
    \item $C_{\text{arête}}$ : Coût unitaire d'une plaquette
\end{itemize}
\end{tcolorbox}

\subsubsection{Efficacité}

L'efficacité mesure la proportion de temps VA dans le temps total de fabrication.

\begin{tcolorbox}[colback=vertclair, colframe=green!50!black, title=Calcul de l'efficacité]
\[
\eta = \frac{T_{VA}}{T_{VA} + T_{\text{orientation}} + T_{\text{accessoire}} + T_{\text{outil}} + T_{\text{plaquette}}}
\]
\end{tcolorbox}

\subsubsection{Indice de risque technique}

L'indice de risque quantifie la probabilité de rupture d'outil ou de défaillance.

\begin{tcolorbox}[colback=orangeclair, colframe=orange!70!black, title=Calcul de l'indice de risque]
\[
IR_{\text{op}} = F_z \times Z \times Ap \times Ae \times \frac{L^3}{D^3}
\]
\[
IR_{\text{gamme}} = 
\begin{cases}
\min_n(IR_{\text{op}}) & \text{(attitude optimiste)} \\
\displaystyle \frac{1}{N}\sum_{n=1}^N IR_{\text{op}} & \text{(attitude neutre)} \\
\max_n(IR_{\text{op}}) & \text{(attitude pessimiste)} \\
-10\log(S^2 + \overline{Y}^{-2}) & \text{(attitude robuste)}
\end{cases}
\]
\end{tcolorbox}

\subsection{Approche Algorithmique pour la Classe C (Algorithmes génétiques)}

\subsubsection{Principe de l'algorithme génétique}

L'algorithme génétique est une métaheuristique inspirée de la théorie de l'évolution. Il permet d'explorer un grand nombre de solutions pour trouver la gamme optimale.

\begin{tcolorbox}[colback=blueclair, colframe=blue!70!black, title=Analogies de l'algorithme génétique]
\begin{tabular}{|l|l|}
\hline
\textbf{Génétique} & \textbf{Application} \\
\hline
Population & Ensemble de gammes \\
Chromosome & Une gamme \\
Gène & Une opération \\
Génération & Itération \\
Fitness & Performance de la gamme \\
Sélection & Conservation des meilleures gammes \\
Croisement & Combinaison de deux gammes \\
Mutation & Modification aléatoire d'une gamme \\
\hline
\end{tabular}
\end{tcolorbox}

\subsubsection{Structure de l'algorithme}

\begin{enumerate}
    \item \textbf{Création de la population initiale}
    \begin{itemize}
        \item Sélection aléatoire d'une gamme enveloppe par zone
        \item Séquencement zone par zone
        \item Mixage aléatoire pour 25\% de la population
    \end{itemize}
    
    \item \textbf{Évaluation des individus}
    \begin{itemize}
        \item Calcul des indicateurs de performance (VA, NVA, coût outil, efficacité)
        \item Normalisation des valeurs
        \item Calcul du macro-indicateur par somme pondérée
    \end{itemize}
    
    \item \textbf{Sélection}
    \begin{itemize}
        \item Conservation des 10\% meilleurs individus
        \item Sélection par tournoi (10\% de la population)
    \end{itemize}
    
    \item \textbf{Croisement}
    \begin{itemize}
        \item Probabilité de croisement $P_c = 0,6$
        \item Échange de gammes enveloppes entre parents
    \end{itemize}
    
    \item \textbf{Mutation}
    \begin{itemize}
        \item Mutation de séquencement : inversion de deux opérations
        \item Mutation de zone : changement d'une gamme enveloppe
        \item Mutation d'outil : remplacement d'un outil par un compatible
    \end{itemize}
    
    \item \textbf{Itération} : Répétition jusqu'à 300 générations
\end{enumerate}

\subsubsection{Paramètres de l'algorithme}

\begin{table}[h]
\centering
\begin{tabular}{|l|c|}
\hline
\textbf{Paramètre} & \textbf{Valeur} \\
\hline
Nombre d'itérations & 300 \\
Taille de la population & 100 \\
Probabilité de croisement & 0,6 \\
Probabilité mutation séquencement & 0,1 \\
Probabilité mutation zone & 0,1 \\
Probabilité mutation outil & 0,6 \\
Meilleurs individus conservés & 10\% \\
Individus gardés par tournoi & 10\% \\
\hline
\end{tabular}
\end{table}

\subsubsection{Exemple de résultats d'optimisation}

\begin{tcolorbox}[colback=grisclair, colframe=gray!70!black, title=Résultats pour une pièce aéronautique]
\begin{tabular}{|l|c|c|c|}
\hline
\textbf{Indicateur} & \textbf{Gamme initiale} & \textbf{Gamme optimisée} & \textbf{Gain} \\
\hline
Temps VA & 2650 min & 1806 min & 32\% \\
Temps NVA & 834 min & 208 min & 75\% \\
Efficacité & 76\% & 90\% & +14\% \\
Coût outil & 7461 € & 2568 € & 66\% \\
Indice risque pessimiste & 13,4 & 6,6 & 51\% \\
\hline
\end{tabular}
\end{tcolorbox}

\subsection{Approche par Simulation pour les Classes A et B}

\subsubsection{Principe de la simulation}

Pour les classes A et B, l'optimisation repose sur la simulation de la production pour valider les performances des gammes.

\subsubsection{Simulation de flux (Classe A)}

\begin{tcolorbox}[colback=vertclair, colframe=green!50!black, title=Objectifs de la simulation de flux]
\begin{itemize}
    \item Valider l'équilibrage des lignes
    \item Identifier les goulets d'étranglement
    \item Optimiser les tailles de lots
    \item Évaluer les temps de cycle
    \item Calculer les taux d'utilisation
\end{itemize}
\end{tcolorbox}

\subsubsection{Équilibrage de ligne}

L'équilibrage de ligne consiste à répartir la charge de travail entre les postes pour minimiser le temps de cycle.

\begin{tcolorbox}[colback=blueclair, colframe=blue!70!black, title=Calcul du taux d'équilibrage]
\[
TE = \frac{\sum_{i=1}^{N} T_i}{N \times \max(T_i)}
\]
Avec :
\begin{itemize}
    \item $TE$ : Taux d'équilibrage (objectif > 85\%)
    \item $T_i$ : Temps de cycle du poste $i$
    \item $N$ : Nombre de postes
\end{itemize}
\end{tcolorbox}

\subsubsection{Simulation de séquencement (Classe B)}

Pour les gammes à variantes, la simulation permet d'optimiser le séquencement des ordres.

\begin{table}[h]
\centering
\begin{tabular}{|p{0.35\textwidth}|p{0.6\textwidth}|}
\hline
\textbf{Règle de séquencement} & \textbf{Objectif} \\
\hline
Mise en cluster & Regrouper les variantes avec mêmes caractéristiques \\
Blocage & Éviter les successions problématiques \\
Proportionnalité & Respecter le mix produit \\
Priorité client & Traiter en priorité les commandes urgentes \\
\hline
\end{tabular}
\end{table}

\subsubsection{Outils de simulation ERP}

Les ERP modernes intègrent des modules de simulation permettant de tester différents scénarios.

\begin{tcolorbox}[colback=orangeclair, colframe=orange!70!black, title=Fonctionnalités de simulation ERP]
\begin{description}
    \item[Planification sur capacité infinie] Simulation sans contrainte de capacité
    \item[Planification sur capacité finie] Prise en compte des contraintes réelles
    \item[Simulation de séquencement] Optimisation de l'ordre de passage
    \item[Analyse des goulots] Identification des ressources critiques
    \item[Scénarios comparatifs] Évaluation de plusieurs options
\end{description}
\end{tcolorbox}

\subsubsection{Indicateurs de performance simulés}

\begin{table}[h]
\centering
\begin{tabular}{|l|c|c|}
\hline
\textbf{Indicateur} & \textbf{Classe A} & \textbf{Classe B} \\
\hline
Taux de rendement synthétique (TRS) & $\geq 85\%$ & $\geq 80\%$ \\
Temps de cycle & Minimal & Variable \\
Temps de changement de série & $\leq 10$ min & $\leq 15$ min \\
Taux de service & $\geq 98\%$ & $\geq 95\%$ \\
Encours moyen & Minimal & Maîtrisé \\
\hline
\end{tabular}
\end{table}

\subsubsection{Exemple de simulation pour une ligne d'assemblage mixte}

\begin{tcolorbox}[colback=grisclair, colframe=gray!70!black, title=Simulation de séquencement - Ligne automobile]
\begin{tabular}{|c|c|c|c|c|}
\hline
\textbf{Position} & \textbf{Variante} & \textbf{Temps cycle (s)} & \textbf{Changement outil} & \textbf{Changement couleur} \\
\hline
1 & Berline & 120 & -- & -- \\
2 & Berline & 120 & Faible & Faible \\
3 & Break & 125 & Faible & Faible \\
4 & Break & 125 & Faible & Faible \\
5 & Berline & 120 & Faible & Élevé \\
6 & Coupé & 130 & Élevé & Moyen \\
7 & Coupé & 130 & Faible & Faible \\
\hline
\end{tabular}
\vspace{0.5em}
\textit{Résultats : Taux d'utilisation = 87\%, Temps de changement total = 28 min, Production = 45 véhicules/jour}
\end{tcolorbox}

\subsubsection{Optimisation itérative}

\begin{enumerate}
    \item \textbf{Définition du scénario de référence} : Gamme initiale, paramètres de base
    \item \textbf{Simulation du scénario} : Calcul des indicateurs de performance
    \item \textbf{Analyse des résultats} : Identification des points d'amélioration
    \item \textbf{Proposition de modifications} : Ajustements des paramètres de gamme
    \item \textbf{Simulation du nouveau scénario} : Évaluation des gains
    \item \textbf{Validation} : Sélection de la meilleure solution
\end{enumerate}

\begin{tcolorbox}[colback=blueclair, colframe=blue!70!black, title=Synthèse : Optimisation des Gammes]
\begin{enumerate}
    \item \textbf{Indicateurs de performance} : VA, NVA, coût outil, efficacité, risque
    \item \textbf{Classe C (unitaire)} : Algorithme génétique pour explorer l'espace des solutions
    \item \textbf{Classe A (grande série)} : Simulation de flux pour équilibrer et optimiser
    \item \textbf{Classe B (variantes)} : Simulation de séquencement pour le mix produit
    \item \textbf{Objectif} : Réduction des coûts et délais, maîtrise des risques
\end{enumerate}
\vspace{0.5em}
\textit{L'optimisation des gammes est un processus itératif combinant analyse des indicateurs, simulation et algorithmes d'optimisation.}
\end{tcolorbox}
% ============================================
% CHAPITRE 9 : INTÉGRATION ET FLUX DE DONNÉES
% ============================================

\section{Intégration et Flux de Données}

\subsection{Du Plan Industriel et Commercial (PIC) au Programme Directeur (PDP)}

Le Plan Industriel et Commercial (PIC) constitue le niveau de planification le plus agrégé, traduisant la stratégie commerciale en objectifs de production.

\subsubsection{Le Plan Industriel et Commercial (PIC)}

\begin{tcolorbox}[colback=blueclair, colframe=blue!70!black, title=Définition du PIC]
Le PIC est un plan de production agrégé exprimé en familles de produits ou en unités de capacité. Il couvre un horizon de 12 à 18 mois et est révisé mensuellement.
\end{tcolorbox}

\begin{table}[h]
\centering
\begin{tabular}{|p{0.3\textwidth}|p{0.65\textwidth}|}
\hline
\textbf{Données d'entrée} & \textbf{Source} \\
\hline
Prévisions de vente & Marketing / Commercial \\
Commandes fermes & Clients \\
Capacités disponibles & Production \\
Stratégies de production & Direction industrielle \\
Objectifs de stock & Supply Chain \\
\hline
\end{tabular}
\end{table}

\subsubsection{Les stratégies de planification du PIC}

\begin{tcolorbox}[colback=grisclair, colframe=gray!70!black, title=Stratégies de planification]
\begin{description}
    \item[Stratégie de poursuite] La production suit la demande. Avantage : faible stock. Inconvénient : variations de charge.
    \item[Stratégie de niveau (lissage)] La production est constante. Avantage : charge stable. Inconvénient : stock variable.
    \item[Stratégie mixte] Combinaison des deux approches avec gestion des stocks saisonniers.
\end{description}
\end{tcolorbox}

\subsubsection{Du PIC au PDP}

Le Programme Directeur de Production (PDP) décline le PIC en un plan détaillé par article fini.

\begin{tcolorbox}[colback=vertclair, colframe=green!50!black, title=Articulation PIC - PDP]
\begin{center}
\textbf{PIC} (Familles de produits, horizons longs)
$\downarrow$ 
\textbf{PDP} (Articles finis, horizon moyen)
$\downarrow$ 
\textbf{MRP} (Composants, horizon court)
\end{center}
\end{tcolorbox}

\subsubsection{Calcul du PDP}

\[
PDP_t = \text{Demande}_t + \text{Stock de sécurité}_t - \text{Stock initial}_{t-1} - \text{En-cours}_t
\]

\begin{table}[h]
\centering
\begin{tabular}{|c|c|c|c|c|c|}
\hline
\textbf{Période} & \textbf{Demande} & \textbf{Stock initial} & \textbf{En-cours} & \textbf{SS} & \textbf{PDP} \\
\hline
Semaine 1 & 100 & 50 & 30 & 20 & 40 \\
Semaine 2 & 120 & 40 & 40 & 20 & 60 \\
Semaine 3 & 150 & 40 & 30 & 20 & 100 \\
Semaine 4 & 130 & 50 & 50 & 20 & 50 \\
\hline
\end{tabular}
\caption{Exemple de calcul PDP}
\end{table}

\subsection{Du PDP aux Ordres de Fabrication (OF) suggérés}

\subsubsection{Planification des Besoins en Composants (MRP)}

La MRP (Material Requirements Planning) transforme le PDP en besoins nets pour les composants.

\begin{tcolorbox}[colback=blueclair, colframe=blue!70!black, title=Principe du calcul MRP]
\begin{enumerate}
    \item \textbf{Besoins bruts} : Besoins issus du PDP et des nomenclatures
    \item \textbf{Besoins nets} : Besoins bruts - stock disponible - en-cours
    \item \textbf{Proposition d'ordres} : Lancement des ordres selon délais d'obtention
    \item \textbf{Échéancier} : Dates de début et de fin des ordres
\end{enumerate}
\end{tcolorbox}

\subsubsection{Algorithme de calcul des besoins nets}

\[
\text{Besoin net}_t = \max(0, \text{Besoin brut}_t - \text{Stock}_{t-1} - \text{En-cours}_t)
\]

\[
\text{Lancement}_t = \text{Besoin net}_{t + \text{Délai}}
\]

\subsubsection{Règles de regroupement (lotissement)}

\begin{table}[h]
\centering
\begin{tabular}{|p{0.35\textwidth}|p{0.6\textwidth}|}
\hline
\textbf{Règle} & \textbf{Description} \\
\hline
Lot pour lot & Lancement exact du besoin \\
Quantité économique & Formule de Wilson \\
Période fixe & Regroupement des besoins sur une période donnée \\
Période variable & Regroupement jusqu'à un coût seuil \\
\hline
\end{tabular}
\end{table}

\subsubsection{Transformation en Ordres de Fabrication}

\begin{enumerate}
    \item \textbf{Ordres planifiés} : Suggestions de l'ERP issues du calcul MRP
    \item \textbf{Transfert en ordres suggérés} : Vérification de cohérence
    \item \textbf{Conversion en ordres fermes} : Après validation des capacités
    \item \textbf{Lancement} : Création des ordres dans le module de pilotage
\end{enumerate}

\subsection{Jalonnements, affermissement et lancement des OF}

\subsubsection{Jalonnement des ordres de fabrication}

Le jalonnement consiste à positionner les ordres dans le temps en fonction des ressources disponibles.

\begin{tcolorbox}[colback=grisclair, colframe=gray!70!black, title=Étapes du jalonnement]
\begin{enumerate}
    \item \textbf{Jalonnement à capacité infinie} : Planification sans contrainte de charge
    \item \textbf{Analyse des charges} : Confrontation capacité disponible / capacité requise
    \item \textbf{Lissage des charges} : Déplacement des ordres pour résoudre les surcharges
    \item \textbf{Jalonnement à capacité finie} : Planification avec contraintes réelles
\end{enumerate}
\end{tcolorbox}

\subsubsection{Affermissement des ordres}

L'affermissement est l'étape qui transforme un ordre planifié en un ordre ferme, engageant les ressources et les matières.

\begin{table}[h]
\centering
\begin{tabular}{|p{0.35\textwidth}|p{0.6\textwidth}|}
\hline
\textbf{Statut} & \textbf{Description} \\
\hline
Planifié & Proposition issue du calcul MRP, modifiable librement \\
Affermi & Vérifié par le planificateur, réservations matières créées \\
Lancé & Transféré à l'atelier, consommations possibles \\
Actif & Travail commencé, encours valorisé \\
Terminé & Fabrication achevée, produit réceptionné \\
Clôturé & Écarts calculés, ordre clos \\
\hline
\end{tabular}
\end{table}

\subsubsection{Lancement des ordres de fabrication}

\begin{tcolorbox}[colback=vertclair, colframe=green!50!black, title=Contrôles avant lancement]
\begin{itemize}
    \item Vérification disponibilité des matières (réservations)
    \item Vérification disponibilité des outils
    \item Vérification capacité des postes de charge
    \item Vérification des qualifications opérateurs
    \item Édition des documents de fabrication (dossier, bons de travail)
    \item Création des ordres magasin pour sortie matières
    \item Génération des ordres de contrôle qualité
\end{itemize}
\end{tcolorbox}

\subsubsection{Les différents statuts de l'ordre de fabrication}

\begin{table}[h]
\centering
\begin{tabular}{|l|l|}
\hline
\textbf{Statut} & \textbf{Actions possibles} \\
\hline
Créé & Modification libre, pas de sortie matière \\
Imprimé & Documents édités, lancement possible \\
Lancé & Sortie matière possible, imputation heures \\
Actif & Première transaction, coûts gelés \\
Fabrication terminée & Pièce achevée, réception possible \\
Terminé & Produits réceptionnés, encore modifiable \\
Fermé & Écarts calculés, plus de modifications \\
Archivé & Données historisées \\
\hline
\end{tabular}
\end{table}

\subsection{Retour d'expérience et boucle d'amélioration continue}

\subsubsection{Calcul des coûts réels}

Le calcul des coûts réels permet de comparer les performances théoriques et réelles.

\begin{tcolorbox}[colback=blueclair, colframe=blue!70!black, title=Composantes des coûts réels]
\[
C_{\text{réel}} = \sum (Q_{\text{matières}} \times P_{\text{matières}}) + \sum (T_{\text{heures}} \times T_{\text{horaire}}) + C_{\text{outils}}
\]
\begin{itemize}
    \item $Q_{\text{matières}}$ : Quantités réellement consommées
    \item $P_{\text{matières}}$ : Prix réels d'achat
    \item $T_{\text{heures}}$ : Temps réellement passés
    \item $T_{\text{horaire}}$ : Taux horaires réels
\end{itemize}
\end{tcolorbox}

\subsubsection{Analyse des écarts}

\begin{table}[h]
\centering
\begin{tabular}{|p{0.35\textwidth}|p{0.6\textwidth}|}
\hline
\textbf{Type d'écart} & \textbf{Description} \\
\hline
Écart de matière & Différence entre quantités standards et réelles \\
Écart de prix matière & Différence entre prix standard et réel \\
Écart de temps & Différence entre temps standards et réels \\
Écart de taux & Différence entre taux horaires standard et réel \\
Écart de rendement & Impact des rebuts et pertes \\
\hline
\end{tabular}
\end{table}

\subsubsection{Boucle d'amélioration continue}

\begin{tcolorbox}[colback=orangeclair, colframe=orange!70!black, title=Cycle PDCA d'amélioration des gammes]
\begin{enumerate}
    \item \textbf{Plan (Planifier)} : Définir les objectifs d'amélioration
    \item \textbf{Do (Réaliser)} : Mettre en œuvre les modifications
    \item \textbf{Check (Vérifier)} : Mesurer les écarts et les gains
    \item \textbf{Act (Agir)} : Standardiser les bonnes pratiques
\end{enumerate}
\end{tcolorbox}

\subsubsection{Capitalisation des retours d'expérience}

\begin{itemize}
    \item \textbf{Temps réels vs temps standards} : Mise à jour des gammes
    \item \textbf{Conditions de coupe optimisées} : Ajustement des paramètres
    \item \textbf{Outils alternatifs} : Enrichissement de la base de données
    \item \textbf{Non-conformités} : Identification des opérations à risque
    \item \textbf{Meilleures pratiques} : Standardisation des solutions efficaces
\end{itemize}

\subsubsection{Indicateurs de pilotage}

\begin{table}[h]
\centering
\begin{tabular}{|l|c|c|}
\hline
\textbf{Indicateur} & \textbf{Objectif} & \textbf{Suivi} \\
\hline
Taux de conformité des gammes & 100\% & Mensuel \\
Écart temps réel / standard & $< 5\%$ & Hebdomadaire \\
Taux de réutilisation des gammes & $> 80\%$ & Trimestriel \\
Délai de mise à jour des gammes & $< 48h$ & Hebdomadaire \\
Nombre de modifications techniques & Variable & Mensuel \\
\hline
\end{tabular}
\end{table}

\subsubsection{Flux d'information complet}

\begin{tcolorbox}[colback=grisclair, colframe=gray!70!black, title=Synthèse du flux de données]
\begin{center}
\textbf{PIC} (Familles, 12-18 mois)\\
$\downarrow$\\
\textbf{PDP} (Articles finis, 3-6 mois)\\
$\downarrow$\\
\textbf{MRP} (Composants, délais d'obtention)\\
$\downarrow$\\
\textbf{Ordres planifiés} (Suggestions)\\
$\downarrow$\\
\textbf{Jalonnement et affermissement} (Validation capacité)\\
$\downarrow$\\
\textbf{Lancement} (Ordres fermes)\\
$\downarrow$\\
\textbf{Exécution} (Production réelle)\\
$\downarrow$\\
\textbf{Calcul des coûts réels} (Analyse des écarts)\\
$\downarrow$\\
\textbf{Mise à jour des gammes} (Amélioration continue)
\end{center}
\end{tcolorbox}

\begin{tcolorbox}[colback=blueclair, colframe=blue!70!black, title=Synthèse : Intégration et Flux de Données]
\begin{enumerate}
    \item \textbf{PIC $\rightarrow$ PDP} : Déclinaison des prévisions en programme directeur
    \item \textbf{PDP $\rightarrow$ MRP} : Calcul des besoins nets pour tous les composants
    \item \textbf{MRP $\rightarrow$ OF} : Transformation des suggestions en ordres fermes
    \item \textbf{Jalonnement} : Planification des ordres avec contraintes de capacité
    \item \textbf{Lancement} : Engagement des ressources et édition des documents
    \item \textbf{Retour d'expérience} : Analyse des écarts et amélioration continue
\end{enumerate}
\vspace{0.5em}
\textit{L'intégration des flux de données garantit la cohérence entre la planification et l'exécution, permettant une amélioration continue des gammes.}
\end{tcolorbox}
% ============================================
% CONCLUSION : VERS UNE GESTION AGILE DES GAMMES
% ============================================

\section*{Conclusion : Vers une Gestion Agile des Gammes}

\subsection*{Récapitulatif des approches par classe}

La gestion des gammes de production ne saurait se réduire à une approche unique. La diversité des contextes industriels impose une méthodologie adaptée à chaque classe de produits.

\begin{tcolorbox}[colback=blueclair, colframe=blue!70!black, title=Synthèse des trois classes]
\begin{center}
\begin{tabular}{|p{0.25\textwidth}|p{0.25\textwidth}|p{0.25\textwidth}|}
\hline
\textbf{Classe A} & \textbf{Classe B} & \textbf{Classe C} \\
\textbf{Grande série} & \textbf{Moyenne série} & \textbf{Unitaire} \\
\hline
Standardisation & Modularité & Flexibilité \\
Optimisation flux & Mass Customization & Gestion de projet \\
PDP / MRP & Configurateur & Planification réseau \\
Post-consommation & Séquencement mixte & Gestion des risques \\
TRS $>$ 85\% & Taux de service $>$ 95\% & Délai projet 100\% \\
\hline
\end{tabular}
\end{center}
\end{tcolorbox}

\subsubsection{Classe A : L'excellence opérationnelle}

Pour les produits de grande série, la performance repose sur :
\begin{itemize}
    \item \textbf{Standardisation} : Nomenclatures et gammes figées, réutilisables
    \item \textbf{Optimisation des flux} : Équilibrage des lignes, réduction des temps morts
    \item \textbf{Post-consommation} : Simplification des transactions de stock
    \item \textbf{Pilotage par indicateurs} : TRS, temps de cycle, taux de défauts
\end{itemize}

\subsubsection{Classe B : L'agilité maîtrisée}

Pour les produits à variantes, la performance repose sur :
\begin{itemize}
    \item \textbf{Modularité} : Conception par modules standardisés
    \item \textbf{Configuration} : Passage de l'article générique à la variante
    \item \textbf{Séquencement mixte} : Production de variantes différentes sur une même ligne
    \item \textbf{Équilibre} : Concilier diversité et productivité
\end{itemize}

\subsubsection{Classe C : La fiabilité projet}

Pour les produits unitaires, la performance repose sur :
\begin{itemize}
    \item \textbf{Approche projet} : Structure PCS avec jalons et activités
    \item \textbf{Gestion des risques} : Identification, modélisation, mitigation
    \item \textbf{Planification réseau} : Chemin critique, marges
    \item \textbf{Traçabilité totale} : Suivi matière, opérations, contrôles
\end{itemize}

\subsection*{Importance de la capitalisation du savoir-faire et de la collaboration}

\subsubsection{La capitalisation du savoir-faire}

La connaissance des méthodistes est un actif stratégique de l'entreprise. Sa capitalisation est essentielle pour :
\begin{itemize}
    \item \textbf{Préserver l'expertise} face au turn-over
    \item \textbf{Accélérer l'intégration} des nouveaux collaborateurs
    \item \textbf{Réutiliser les bonnes pratiques} d'un projet à l'autre
    \item \textbf{Éviter les erreurs déjà rencontrées}
\end{itemize}

\begin{tcolorbox}[colback=grisclair, colframe=gray!70!black, title=Outils de capitalisation du savoir-faire]
\begin{description}
    \item[Base de données outils] Références, conditions de coupe, retours d'expérience
    \item[Bibliothèque de gammes] Gammes standards par famille de produits
    \item[Fiches de retour d'expérience] Analyse des écarts et des non-conformités
    \item[Matrice des compétences] Qualifications et expertises des opérateurs
    \item[Base de connaissances] Solutions techniques, astuces, alertes
\end{description}
\end{tcolorbox}

\subsubsection{La collaboration entre métiers}

La performance des gammes repose sur une collaboration étroite entre les différents acteurs :

\begin{tcolorbox}[colback=vertclair, colframe=green!50!black, title=Collaboration conception - méthodes - production]
\begin{center}
\begin{tabular}{c}
\textbf{Bureau d'Études} (Conception)\\
$\updownarrow$ \\
\textbf{Bureau des Méthodes} (Industrialisation)\\
$\updownarrow$ \\
\textbf{Production} (Exécution)
\end{tabular}
\end{center}
\vspace{0.5em}
\textit{La conception doit intégrer les contraintes de fabrication (DFM - Design For Manufacturing).}
\textit{Les méthodes doivent transmettre des gammes exploitables en production.}
\textit{La production doit remonter les écarts pour améliorer les gammes.}
\end{tcolorbox}

\subsection*{Perspectives : Intégration des données temps réel (MES) pour des gammes adaptatives}

\subsubsection{Du MES à la gamme adaptative}

Les Manufacturing Execution Systems (MES) collectent en temps réel les données de production. Leur intégration ouvre la voie à des gammes adaptatives.

\begin{table}[h]
\centering
\begin{tabular}{|p{0.35\textwidth}|p{0.6\textwidth}|}
\hline
\textbf{Données MES} & \textbf{Application à la gamme} \\
\hline
Temps réels d'opération & Ajustement des temps standards \\
Consommations matières réelles & Mise à jour des coefficients de rebut \\
Usure outil mesurée & Optimisation des durées de vie \\
Vibrations / efforts de coupe & Ajustement des conditions de coupe \\
Taux de non-conformité & Révision des opérations critiques \\
Disponibilité machine & Replanification dynamique \\
\hline
\end{tabular}
\end{table}

\subsubsection{La gamme adaptative : un concept émergent}

\begin{tcolorbox}[colback=orangeclair, colframe=orange!70!black, title=Vers la gamme 4.0]
\begin{description}
    \item[Gamme statique] Définie a priori, figée jusqu'à modification
    \item[Gamme réactive] Ajustée après analyse des écarts
    \item[Gamme adaptative] S'ajuste en temps réel aux conditions de production
\end{description}
\end{tcolorbox}

\subsubsection{Les briques technologiques de la gamme adaptative}

\begin{enumerate}
    \item \textbf{Capteurs connectés} : Mesure en continu des conditions de coupe
    \item \textbf{Communication MES-ERP} : Remontée automatique des données
    \item \textbf{Algorithmes d'optimisation} : Recalcul des paramètres en temps réel
    \item \textbf{Intelligence artificielle} : Apprentissage et prédiction
    \item \textbf{Jumeau numérique} : Simulation avant exécution
\end{enumerate}

\subsubsection{Exemple : Optimisation temps réel des conditions de coupe}

\begin{tcolorbox}[colback=blueclair, colframe=blue!70!black, title=Scénario d'optimisation adaptative]
\begin{enumerate}
    \item \textbf{Capteur} : Mesure l'usure de l'outil en temps réel
    \item \textbf{Analyse} : L'IA prédit la durée de vie restante
    \item \textbf{Décision} : Ajustement automatique de l'avance
    \item \textbf{Action} : Modification des paramètres CN en ligne
    \item \textbf{Capitalisation} : Enrichissement de la base de connaissances
\end{enumerate}
\textit{Résultat : Durée de vie outil augmentée de 15\%, qualité maintenue}
\end{tcolorbox}

\subsubsection{Bénéfices attendus}

\begin{table}[h]
\centering
\begin{tabular}{|l|c|}
\hline
\textbf{Bénéfice} & \textbf{Impact potentiel} \\
\hline
Réduction des temps d'arrêt & $-20\%$ \\
Augmentation de la durée de vie outil & $+15\%$ \\
Réduction des rebuts & $-30\%$ \\
Amélioration du TRS & $+10\%$ \\
Réactivité face aux aléas & $+50\%$ \\
\hline
\end{tabular}
\end{table}

\subsubsection{Feuille de route vers la gamme adaptative}

\begin{tcolorbox}[colback=grisclair, colframe=gray!70!black, title=Plan d'évolution]
\begin{enumerate}
    \item \textbf{Niveau 1} : Collecte des données MES, analyse a posteriori
    \item \textbf{Niveau 2} : Remontée automatique vers l'ERP, révision périodique
    \item \textbf{Niveau 3} : Alertes en temps réel, ajustements manuels
    \item \textbf{Niveau 4} : Optimisation automatique, supervision humaine
    \item \textbf{Niveau 5} : Gamme autonome et prédictive
\end{enumerate}
\end{tcolorbox}

% ============================================
% CONCLUSION FINALE
% ============================================

\begin{tcolorbox}[colback=blueclair, colframe=blue!70!black, title=Conclusion Générale]
\begin{center}
\textbf{La gamme de production : un levier stratégique de performance}
\end{center}
\vspace{0.5em}

La maîtrise des gammes de production est aujourd'hui un facteur clé de compétitivité. Ce document a démontré qu'une approche unique ne peut répondre à la diversité des contextes industriels.

\vspace{0.5em}
\textbf{Trois classes de produits, trois logiques de gestion :}
\begin{itemize}
    \item \textbf{Classe A} : Optimisation des flux, standardisation, répétitivité
    \item \textbf{Classe B} : Modularité, configuration, séquencement mixte
    \item \textbf{Classe C} : Flexibilité, gestion de projet, maîtrise des risques
\end{itemize}

\vspace{0.5em}
\textbf{Une intégration réussie repose sur :}
\begin{itemize}
    \item La \textbf{formalisation rigoureuse} des données dans l'ERP
    \item L'\textbf{optimisation} par simulation et algorithmes
    \item La \textbf{capitalisation} du savoir-faire méthodiste
    \item La \textbf{collaboration} entre conception, méthodes et production
\end{itemize}

\vspace{0.5em}
\textbf{Perspectives :}
L'intégration des données temps réel via le MES ouvre la voie vers des gammes adaptatives, capables de s'ajuster en continu aux conditions réelles de production. L'intelligence artificielle et le jumeau numérique permettront à terme des gammes prédictives, anticipant les aléas et optimisant en continu la performance.

\vspace{0.5em}
\textit{La gamme de production n'est plus un document figé, mais un outil vivant au cœur du système d'information industriel.}
\end{tcolorbox}

\begin{center}
\textbf{Fin du document}
\end{center}

\end{document}

